\frfilename{mt244.tex}
\versiondate{6.3.09}
\copyrightdate{1995}

\def\chaptername{Ruang fungsi}
\def\sectionname{$L^{p}$}

\newsection{244}

Melanjutkan peninjauan kita atas ruang-ruang Banach klasik, kita sampai
pada ruang-ruang $L^p$ untuk $1<p<\infty$. Kasus $p=2$ lebih penting
daripada gabungan semua kasus lainnya, dan masuk akal, bahkan mungkin
dianjurkan, untuk pertama-tama membaca bagian ini hanya dengan memikirkan
kasus tersebut. Namun ruang-ruang lainnya memberikan contoh yang
mendidik dan tetap menjadi bagian dasar pendidikan setiap analis
fungsional.

\leader{244A}{Definisi}
Misalkan $(X,\Sigma,\mu)$ sebarang ruang ukur, dan
$p\in\ooint{1,\infty}$. Tuliskan $\eusm L^p=\eusm L^p(\mu)$ untuk
himpunan fungsi $f\in\eusm L^0=\eusm L^0(\mu)$ sedemikian sehingga
$|f|^p$ terintegralkan, dan $L^p=L^p(\mu)$ untuk
$\{f^{\ssbullet}:f\in\eusm L^p\}\subseteq L^0=L^0(\mu)$.

\cmmnt{Perhatikan bahwa jika $f\in\eusm L^p$, $g\in\eusm L^0$ dan
$f\eae g$, maka $|f|^p\eae|g|^p$, sehingga $|g|^p$ terintegralkan dan
$g\in\eusm L^p$; dengan demikian}
$\eusm L^p=\{f:f\in\eusm L^0,\,f^{\ssbullet}\in L^p\}$.

\cmmnt{Sebagai alternatif, kita dapat mendefinisikan $u^p$ setiap kali
$u\in L^0$, $u\ge 0$ dengan menuliskan
$(f^{\ssbullet})^p=(f^p)^{\ssbullet}$ untuk setiap $f\in\eusm L^0$
sedemikian sehingga $f(x)\ge 0$ untuk setiap $x\in\dom f$ (bandingkan
241I), dan mengatakan bahwa
$L^p=\{u:u\in L^0,\,|u|^p\in L^1(\mu)\}$.
}

\leader{244B}{Teorema} Misalkan $(X,\Sigma,\mu)$ sebarang ruang ukur,
dan $p\in[1,\infty]$.

(a) $L^p=L^p(\mu)$ merupakan subruang linear dari $L^0=L^0(\mu)$.

(b) Jika $u\in L^p$, $v\in L^0$ dan $|v|\le |u|$, maka $v\in L^p$.
Akibatnya $|u|$, $u\vee v$ dan $u\wedge v$ termasuk dalam $L^p$ untuk
semua $u$, $v\in L^p$.

\proof{ Kasus $p=1$, $p=\infty$ tercakup oleh 242B, 242C dan 243B;
jadi saya andaikan bahwa $1<p<\infty$.

\medskip

{\bf (a)(i)} Andaikan $f$, $g\in\eusm L^p=\eusm L^p(\mu)$. Jika $a$,
$b\in\Bbb R$, maka $|a+b|^p\le 2^p\max(|a|^p,|b|^p)$, sehingga
$|f+g|^p\leae 2^p(|f|^p\vee|g|^p)$; sekarang
$|f+g|^p\in\eusm L^0$ dan
$2^p(|f|^p\vee|g|^p)\in\eusm L^1$, sehingga
$|f+g|^p\in\eusm L^1$. Jadi $f+g\in\eusm L^p$ untuk semua $f$,
$g\in \eusm L^p$; segera diperoleh bahwa $u+v\in L^p$ untuk semua $u$,
$v\in L^p$.

\medskip
\quad{\bf (ii)} Jika $f\in \eusm L^p$ dan $c\in\Bbb R$, maka
$|cf|^p=|c|^p|f|^p\in\eusm L^1$, sehingga $cf\in\eusm L^p$.
Dengan demikian $cu\in L^p$ setiap kali $u\in L^p$ dan $c\in\Bbb R$.

\medskip

{\bf (b)(i)} Nyatakan $u$ sebagai $f^{\ssbullet}$ dan $v$ sebagai
$g^{\ssbullet}$, dengan $f\in\eusm L^p$ dan $g\in\eusm L^0$. Maka
$|g|\leae|f|$, sehingga $|g|^p\leae|f|^p$ dan $|g|^p$ terintegralkan;
dengan demikian $g\in\eusm L^p$ dan $v\in L^p$.

\medskip

\quad{\bf (ii)} Sekarang $|\,|u|\,|=|u|$, sehingga $|u|\in L^p$ setiap
kali $u\in L^p$. Akhirnya,
$u\vee v=\bover12(u+v+|u-v|)$ dan
$u\wedge v=\bover12(u+v-|u-v|)$ termasuk dalam $L^p$ untuk semua $u$,
$v\in L^p$.
}

\leader{244C}{Struktur terurut $L^p$} Misalkan $(X,\Sigma,\mu)$ sebarang
ruang ukur, dan $p\in[1,\infty]$. Maka\cmmnt{ 244B cukup untuk menjamin
bahwa} urutan parsial yang diwarisi dari $L^0(\mu)$ menjadikan
$L^p(\mu)$ suatu ruang Riesz\cmmnt{ (bandingkan 242C, 243C)}.

\leader{244D}{Norma $L^p$} Misalkan $(X,\Sigma,\mu)$ suatu ruang ukur,
dan $p\in\ooint{1,\infty}$.

\spheader 244Da Untuk $f\in\eusm L^p=\eusm L^p(\mu)$, tetapkan
$\|f\|_p=(\int|f|^p)^{1/p}$.
Jika $f$, $g\in\eusm L^p$ dan $f\eae g$, maka
$|f|^p\eae|g|^p$, sehingga $\|f\|_p=\|g\|_p$. Dengan demikian kita dapat
mendefinisikan $\|\,\|_p:L^p=L^p(\mu)\to\coint{0,\infty}$ dengan
menuliskan $\|f^{\ssbullet}\|_p=\|f\|_p$ untuk setiap $f\in\eusm L^p$.

\cmmnt{Sebagai alternatif, kita dapat mengatakan saja bahwa
$\|u\|_p=(\int|u|^p)^{1/p}$ untuk setiap $u\in L^p=L^p(\mu)$.}

\spheader 244Db\cmmnt{ Notasi $\|\,\|_p$ membawa janji bahwa ia
merupakan norma pada $L^p$; ini memang benar, tetapi saya menunda
buktinya hingga 244F di bawah. Namun untuk sementara, mari kita perhatikan
saja bahwa} $\|cu\|_p=|c|\|u\|_p$ untuk semua $u\in L^p$ dan
$c\in\Bbb R$, dan\cmmnt{ bahwa} jika $\|u\|_p=0$, maka\cmmnt{
$\int|u|^p=0$, sehingga $|u|^p=0$ dan} $u=0$.

\spheader 244Dc Jika $|u|\le |v|$ dalam $L^p$, maka
$\|u\|_p\le\|v\|_p$\cmmnt{; ini karena $|u|^p\le|v|^p$}.

\leader{244E}{}\cmmnt{ Sekarang saya mengerjakan lema-lema yang
diperlukan untuk menunjukkan bahwa $\|\,\|_p$ merupakan norma pada $L^p$
dan, pada akhirnya, bahwa dual ruang bernorma dari $L^p$ dapat
diidentifikasi dengan $L^q$ yang sesuai.

\medskip

\noindent}{\bf Lema} Andaikan $(X,\Sigma,\mu)$ suatu ruang ukur, dan
$p$, $q\in\ooint{1,\infty}$ sedemikian sehingga
${1\over p}+{1\over q}=1$.

(a) $ab\le{1\over p}a^p+{1\over q}b^q$ untuk semua bilangan real
$a$, $b\ge 0$.

(b)(i) $f\times g$ terintegralkan dan

\Centerline{$|\int f\times g|\le\int|f\times g|\le\|f\|_p\|g\|_q$}

\noindent untuk semua $f\in \eusm L^p=\eusm L^p(\mu)$,
$g\in\eusm L^q=\eusm L^q(\mu)$;

\quad(ii) $u\times v\in L^1=L^1(\mu)$ dan

\Centerline{$|\int u\times v|\le\|u\times v\|_1\le\|u\|_p\|v\|_q$}

\noindent untuk semua $u\in L^p=L^p(\mu)$,
$v\in L^q=L^q(\mu)$.


\proof{{\bf (a)} Jika $a$ atau $b$ bernilai $0$, pernyataan ini sepele.
Jika keduanya tak nol, kita dapat berargumen sebagai berikut. Fungsi
$x\mapsto x^{1/p}:\coint{0,\infty}\to\Bbb R$ cekung, dengan turunan kedua
yang secara ketat kurang dari $0$, sehingga seluruhnya terletak di bawah
setiap garis singgungnya; khususnya, di bawah garis singgungnya pada titik
$(1,1)$, yang mempunyai persamaan $y=1+{1\over p}(x-1)$. Jadi kita
mempunyai

\Centerline{$x^{1/p}\le\Bover1px+1-\Bover1p=\Bover1px+\Bover1q$}

\noindent untuk setiap $x\in\coint{0,\infty}$. Jadi jika $c$, $d>0$,
maka

\Centerline{$(\Bover{c}{d})^{1/p}\le
\Bover1p\Bover{c}{d}+\Bover1q$;}

\noindent dengan mengalikan kedua ruas dengan $d$,

\Centerline{$c^{1/p}d^{1/q}\le\Bover1pc+\Bover1qd$;}

\noindent dengan menetapkan $c=a^p$ dan $d=b^q$, kita memperoleh

\Centerline{$ab\le\Bover1pa^p+\Bover1qb^q$,}

\noindent sebagaimana dinyatakan.

\medskip

{\bf (b)(i)}($\alpha$) Andaikan terlebih dahulu bahwa
$\|f\|_p=\|g\|_q=1$. Untuk setiap $x\in\dom f\cap\dom g$, kita mempunyai

\Centerline{$|f(x)g(x)|\le\Bover1p|f(x)|^p+\Bover1q|g(x)|^q$}

\noindent menurut (a). Jadi

\Centerline{$|f\times g|\leae
\Bover1p|f|^p+\Bover1q|g|^q\in\eusm L^1(\mu)$}

\noindent dan $f\times g$ terintegralkan; selain itu,

\Centerline{$\int|f\times g|\le\Bover1p\int|f|^p+\Bover1q\int|g|^q
=\Bover1p\|f\|_p^p+\Bover1q\|g\|_q^q
=\Bover1p+\Bover1q=1$.}

\noindent ($\beta$) Jika $\|f\|_p=0$, maka $\int|f|^p=0$, sehingga
$|f|^p\eae\tbf{0}$, $f\eae\tbf{0}$, $f\times g\eae\tbf{0}$ dan

\Centerline{$\int|f\times g|=0=\|f\|_p\|g\|_q$.}

\noindent Demikian pula, jika $\|g\|_q=0$, maka $g\eae\tbf{0}$ dan
sekali lagi

\Centerline{$\int|f\times g|=0=\|f\|_p\|g\|_q$.}

\noindent($\gamma$) Akhirnya, untuk $f\in\eusm L^p$,
$g\in\eusm L^q$ umum sedemikian sehingga $c=\|f\|_p$ dan
$d=\|g\|_q$ keduanya tak nol, kita mempunyai
$\|\bover1cf\|_p=\|\bover1dg\|_q=1$, sehingga

\Centerline{$f\times g=cd(\Bover1cf\times\Bover1dg)$}

\noindent terintegralkan, dan

\Centerline{$\int|f\times g|
=cd\int|\Bover1cf\times\Bover1dg|\le cd$,}

\noindent sebagaimana diperlukan.

\medskip

\quad{\bf (ii)} Sekarang jika $u\in L^p$, $v\in L^q$, ambil
$f\in\eusm L^p$, $g\in \eusm L^q$ sedemikian sehingga
$u=f^{\ssbullet}$ dan $v=g^{\ssbullet}$; $f\times g$ terintegralkan,
sehingga $u\times v\in L^1$, dan

\Centerline{$|\int u\times v|\le\|u\times v\|_1
=\int|f\times g|\le\|f\|_p\|g\|_q=\|u\|_p\|v\|_q$.}
}

\cmmnt{\medskip

\noindent{\bf Catatan} Bagian (b) adalah `{\bf pertidaksamaan
H\"older}'. Dalam kasus $p=q=2$, ini adalah `{\bf pertidaksamaan Cauchy}'.
}

\leader{244F}{Proposisi} Misalkan $(X,\Sigma,\mu)$ suatu ruang ukur dan
$p\in\ooint{1,\infty}$. Tetapkan $q=p/(p-1)$, sehingga
${1\over p}+{1\over q}=1$.

(a) Untuk setiap $u\in L^p=L^p(\mu)$,
$\|u\|_p=\max\{\int u\times v:v\in L^q(\mu),\,\|v\|_q\le 1\}$.

(b) $\|\,\|_p$ merupakan norma pada $L^p$.

\proof{{\bf (a)} Untuk $u\in L^p$, tetapkan

\Centerline{$\tau(u)
=\sup\{\int u\times v:v\in L^q(\mu),\,\|v\|_q\le 1\}$.}

\noindent Menurut 244E(b-ii), $\|u\|_p\ge\tau(u)$. Jika
$\|u\|_p=0$, maka tentu

\Centerline{$0=\|u\|_p=\tau(u)
=\max\{\int u\times v:v\in L^q(\mu),\,\|v\|_q\le 1\}$.}

\noindent Jika $\|u\|_p=c>0$, tinjau

\Centerline{$v=c^{-p/q}\sgn u\times|u|^{p/q}$,}

\noindent dengan, untuk $a\in\Bbb R$, saya menuliskan
$\sgn a=|a|/a$ jika $a\ne 0$, dan $0$ jika $a=0$, sehingga
$\sgn:\Bbb R\to\Bbb R$ merupakan fungsi terukur Borel; untuk
$f\in\eusm L^0$, saya menuliskan $(\sgn f)(x)=\sgn(f(x))$ untuk
$x\in\dom f$, sehingga $\sgn f\in\eusm L^0$; dan untuk
$f\in\eusm L^0$, saya menuliskan
$\sgn(f^{\ssbullet})=(\sgn f)^{\ssbullet}$ untuk mendefinisikan suatu
fungsi $\sgn:L^0\to L^0$ (bandingkan 241I). Maka
$v\in L^q=L^q(\mu)$ dan

\Centerline{$\|v\|_q=(\int|v|^q)^{1/q}=c^{-p/q}(\int|u|^p)^{1/q}
=c^{-p/q}c^{p/q}=1$.}

\noindent Jadi

$$\eqalign{\tau(u)\ge\int u\times v
&=c^{-p/q}\int\sgn u\times|u|\times\sgn u\times|u|^{p/q}\cr
&=c^{-p/q}\int|u|^{1+{p\over q}}
=c^{-p/q}\int|u|^p
=c^{p-{p\over q}}
=c,\cr}$$

\noindent dengan mengingat bahwa $1+{p\over q}=p$,
$p-{p\over q}=1$. Jadi $\tau(u)\ge\|u\|_p$ dan

\Centerline{$\tau(u)=\|u\|_p=\int u\times v$.}

\medskip

{\bf (b)} Mengingat catatan dalam 244Db, saya hanya perlu memeriksa bahwa
$\|u+v\|_p\le\|u\|_p+\|v\|_p$ untuk semua $u$, $v\in L^p$. Namun,
untuk $u$ dan $v$ yang diberikan, misalkan $w\in L^q$ sedemikian sehingga
$\|w\|_q=1$ dan $\int (u+v)\times w=\|u+v\|_p$. Maka

\Centerline{$\|u+v\|_p=\int(u+v)\times w
=\int u\times w+\int v\times w\le\|u\|_p+\|v\|_p$,}

\noindent sebagaimana diperlukan.
}%end of proof of 244F

\cmmnt{\medskip

\noindent{\bf Catatan} Pertidaksamaan segitiga
`$\|u+v\|_p\le\|u\|_p+\|v\|_p$' adalah {\bf pertidaksamaan Minkowski}.}

\leader{244G}{Teorema} Misalkan $(X,\Sigma,\mu)$ sebarang ruang ukur,
dan $p\in[1,\infty]$. Maka $L^p=L^p(\mu)$ merupakan kisi Banach di bawah
normanya $\|\,\|_p$.

\proof{ Kasus $p=1$, $p=\infty$ tercakup oleh 242F dan 243E, jadi mari
kita andaikan bahwa $1<p<\infty$. Kita telah mengetahui bahwa
$\|u\|_p\le\|v\|_p$ setiap kali $|u|\le|v|$, sehingga hanya tersisa
untuk menunjukkan bahwa $L^p$ lengkap.

Misalkan $\sequencen{u_n}$ suatu barisan dalam $L^p$ sedemikian sehingga
$\|u_{n+1}-u_n\|_p\le 4^{-n}$ untuk setiap $n\in\Bbb N$. Perhatikan bahwa

\Centerline{$\|u_n\|_p\le\|u_0\|_p+\sum_{k=0}^{n-1}\|u_{k+1}-u_k\|_p
\le\|u_0\|_p+\sum_{k=0}^{\infty}4^{-k}\le\|u_0\|_p+2$}

\noindent untuk setiap $n$. Untuk setiap $n\in\Bbb N$, pilih
$f_n\in\eusm L^p$ sedemikian sehingga $f_0^{\ssbullet}=u_{0}$,
$f_n^{\ssbullet}=u_{n}-u_{n-1}$ untuk $n\ge 1$; lakukan ini sedemikian
sehingga $\dom f_n=X$ dan $f_n$ terukur terhadap $\Sigma$ (241Bk). Maka
$\|f_n\|_p\le 4^{-n+1}$ untuk $n\ge 1$.

Untuk $m$, $n\in\Bbb N$, tetapkan

\Centerline{$E_{mn}=\{x:|f_m(x)|\ge 2^{-n}\}\in\Sigma$.}

\noindent Maka $|f_m(x)|^p\ge 2^{-np}$ untuk $x\in E_{mn}$, sehingga

\Centerline{$2^{-np}\mu E_{mn}\le\int|f_m|^p<\infty$}

\noindent dan $\mu E_{mn}<\infty$.
Jadi $\chi E_{mn}\in\eusm L^q=\eusm L^q(\mu)$ dan

\Centerline{$\int_{E_{mn}}|f_k|=\int|f_k|\times\chi E_{mn}
\le\|f_k\|_p\|\chi E_{mn}\|_q$}

\noindent untuk setiap $k$, menurut 244E(b-i). Dengan demikian

\Centerline{$\sum_{k=0}^{\infty}\int_{E_{mn}}|f_k|
\le\|\chi E_{mn}\|_q\sum_{k=0}^{\infty}\|f_k\|_p<\infty$,}

\noindent dan $\sum_{k=0}^{\infty}f_k(x)$ ada untuk hampir setiap
$x\in E_{mn}$, menurut 242E. Ini benar untuk semua $m$, $n\in\Bbb N$.
Namun jika $x\in X\setminus\bigcup_{m,n\in\Bbb N}E_{mn}$,
$f_n(x)=0$ untuk setiap $n$, sehingga
$\sum_{k=0}^{\infty}f_k(x)$ tentu ada. Jadi
$g(x)=\sum_{k=0}^{\infty}f_k(x)$ terdefinisi dalam $\Bbb R$ untuk hampir
setiap $x\in X$.

Tetapkan $g_n=\sum_{k=0}^nf_k$; maka
$g_n^{\ssbullet}=u_{n}\in L^p$ untuk setiap $n$, dan
$g(x)=\lim_{n\to\infty}g_n(x)$ terdefinisi untuk hampir setiap $x$.
Sekarang tinjau $|g|^p\eae\lim_{n\to\infty}|g_n|^p$. Kita mengetahui
bahwa

\Centerline{$\liminf_{n\to\infty}\int|g_n|^p
=\liminf_{n\to\infty}\|u_{n}\|_p^p\le(2+\|u_{0}\|_p)^p<\infty$,}

\noindent sehingga menurut Lema Fatou

\Centerline{$\int|g|^p\le\liminf_{k\to\infty}\int|g_k|^p<\infty$.}

\noindent Jadi $u=g^{\ssbullet}\in L^p$. Selain itu, untuk sebarang
$m\in\Bbb N$,

$$\eqalign{\int|g-g_m|^p
&\le\liminf_{n\to\infty}\int|g_n-g_m|^p
=\liminf_{n\to\infty}\|u_{n}-u_{m}\|_p^p\cr
&\le\liminf_{n\to\infty}\sum_{k=m}^{n-1}4^{-kp}
=\sum_{k=m}^{\infty}4^{-kp}
=4^{-mp}/(1-4^{-p}).\cr}$$

\noindent Jadi

\Centerline{$\|u-u_{m}\|_p=(\int|g-g_m|^p)^{1/p}
\le 4^{-m}/(1-4^{-p})^{1/p}\to 0$}

\noindent ketika $m\to\infty$. Jadi $u=\lim_{m\to\infty}u_m$ dalam
$L^p$. Karena $\sequencen{u_n}$ sembarang, $L^p$ lengkap.
}

\leader{244H}{}\cmmnt{ Mengikuti 242M dan 242O, saya perhatikan bahwa
$L^p$ berperilaku seperti $L^1$ sehubungan dengan subruang-subruang padat
tertentu.

\medskip

\noindent}{\bf Proposisi} (a) Misalkan $(X,\Sigma,\mu)$ sebarang ruang
ukur, dan $p\in\coint{1,\infty}$. Maka ruang $S$ dari kelas-kelas
ekuivalensi fungsi sederhana-$\mu$ merupakan subruang linear padat dari
$L^p=L^p(\mu)$.

(b) Misalkan $X$ sebarang subhimpunan dari $\BbbR^r$, dengan $r\ge 1$,
dan misalkan $\mu$ ukuran subruang pada $X$ yang diinduksi oleh ukuran
Lebesgue pada $\BbbR^r$. Tuliskan $C_k$ untuk himpunan fungsi kontinu
(terbatas) $g:\BbbR^r\to\Bbb R$ sedemikian sehingga
$\{x:g(x)\ne 0\}$ terbatas, dan $S_0$ untuk ruang kombinasi linear fungsi
berbentuk $\chi I$, dengan $I\subseteq\BbbR^r$ suatu interval setengah
terbuka yang terbatas. Maka $\{(g\restr X)^{\ssbullet}:g\in C_k\}$ dan
$\{(h\restr X)^{\ssbullet}:h\in S_0\}$ padat dalam $L^p(\mu)$.

\proof{{\bf (a)} Saya mengulangi argumen 242M dengan sedikit perubahan.

\medskip

\quad{\bf (i)} Andaikan $u\in L^p(\mu)$, $u\ge 0$ dan $\epsilon>0$.
Nyatakan $u$ sebagai $f^{\ssbullet}$ dengan
$f:X\to\coint{0,\infty}$ suatu fungsi terukur. Misalkan
$g:X\to\Bbb R$ suatu fungsi sederhana sedemikian sehingga
$0\le g\le f^p$ dan $\int g\ge\int f^p-\epsilon^p$. Tetapkan
$h=g^{1/p}$. Maka $h$ suatu fungsi sederhana dan $h\le f$. Karena $p>1$,
$(f-h)^p+h^p\le f^p$ dan

\Centerline{$\int(f-h)^p\le \int f^p-g\le\epsilon^p$,}

\noindent sehingga

\Centerline{$\|u-h^{\ssbullet}\|_p
=(\int|f-h|^p)^{1/p}\le\epsilon$,}

\noindent sedangkan $h^{\ssbullet}\in S$.

\medskip

\quad{\bf (ii)} Untuk $u\in L^p$ umum dan $\epsilon>0$, $u$ dapat
dinyatakan sebagai $u^+-u^-$, dengan $u^+=u\vee 0$,
$u^-=(-u)\vee 0$ termasuk dalam $L^p$ dan nonnegatif. Menurut (i), kita
dapat menemukan $v_1$, $v_2\in S$ sedemikian sehingga
$\|u^+-v_1\|_p\le\bover12\epsilon$ dan
$\|u^--v_2\|_p\le\bover12\epsilon$, sehingga $v=v_1-v_2\in S$ dan
$\|u-v\|_p\le\epsilon$. Karena $u$ dan $\epsilon$ sembarang, $S$ padat.

\medskip

{\bf (b)} Sekali lagi, semua gagasan terdapat dalam 242O; perubahan yang
diperlukan ada pada rumus, bukan pada metode. Perubahan itu akan lebih
mudah jika sejak awal saya mencatat bahwa setiap kali $f_1$,
$f_2\in\eusm L^p(\mu)$ dan $\int|f_1|^p\le\epsilon^p$,
$\int|f_2|^p\le\delta^p$ (dengan $\epsilon$, $\delta\ge 0$), maka
$\int|f_1+f_2|^p\le(\epsilon+\delta)^p$; ini hanyalah pertidaksamaan
segitiga untuk $\|\,\|_p$ (244Fb). Saya juga akan secara teratur
menyatakan hubungan yang dituju dalam bentuk
`$\int_X|f-g|^p\le\epsilon^p$', `$\int_X|f-h|^p\le\epsilon^p$'.
Sekarang mari kita jalankan argumen 242Oa, dengan lebih ringkas daripada
sebelumnya.

\medskip

\quad{\bf (i)} Andaikan terlebih dahulu bahwa $f=\chi I\restr X$ dengan
$I\subseteq\BbbR^r$ suatu interval setengah terbuka yang terbatas.
Seperti sebelumnya, kita dapat menetapkan $h=\chi I$ dan memperoleh
$\int_X|f-h|^p=0$. Kali ini, gunakan konstruksi yang sama untuk menemukan
suatu interval $J$ dan suatu fungsi $g\in C_k$ sedemikian sehingga
$\chi I\le g\le \chi J$ dan
$\mu_r(J\setminus I)\le\epsilon^p$; ini akan menjamin bahwa
$\int_X|f-g|^p\le\epsilon^p$.

\medskip

\quad{\bf (ii)} Sekarang andaikan $f=\chi(X\cap E)$ dengan
$E\subseteq\BbbR^r$ suatu himpunan berukuran hingga. Maka, dengan alasan
yang sama seperti sebelumnya, terdapat keluarga saling lepas
$I_0,\ldots,I_n$ dari interval-interval setengah terbuka sedemikian
sehingga
$\mu_r(E\symmdiff\bigcup_{j\le n}I_j)\le(\bover12\epsilon)^p$.
Dengan demikian $h=\sum_{j=0}^n\chi I_j\in S_0$ dan
$\int_X|f-h|^p\le(\bover12\epsilon)^p$.
Dan (i) memberi tahu kita bahwa untuk setiap $j\le n$ terdapat
$g_j\in C_k$ sedemikian sehingga
$\int_X|g_j-\chi I_j|^p\le(\epsilon/2(n+1))^p$, sehingga
$g=\sum_{j=0}^ng_j\in C_k$ dan $\int_X|f-g|^p\le\epsilon^p$.

\medskip

\quad{\bf (iii)} Peralihan ke fungsi sederhana, lalu ke anggota
sembarang $\eusm L^p(\mu)$, persis seperti sebelumnya, tetapi menggunakan
$\|f\|_p$ sebagai pengganti $\int_X|f|$. Akhirnya, penerjemahan dari
$\eusm L^p$ ke $L^p$ kembali langsung -- dengan mengingat, seperti
sebelumnya, untuk memeriksa bahwa $g\restr X$, $h\restr X$ termasuk dalam
$\eusm L^p(\mu)$ setiap kali $g\in C_k$ dan $h\in S_0$.
}%end of proof of 244H

\leader{*244I}{Korolari} Dalam konteks 244Hb, $L^p(\mu)$ separabel.

\proof{ Misalkan $A$ himpunan

\Centerline{$\{(\sum_{j=0}^nq_j\chi(\coint{a_j,b_j}\cap X))^{\ssbullet}:
n\in\Bbb N$, $q_0,\ldots,q_n\in\Bbb Q$,
$a_0,\ldots,a_n,b_0,\ldots,b_n\in\BbbQ^r\}$.}

\noindent Maka $A$ merupakan subhimpunan terhitung dari $L^p(\mu)$, dan
penutupannya harus memuat
$(\sum_{j=0}^nc_j\chi(\coint{a_j,b_j}\cap X))^{\ssbullet}$ setiap kali
$c_0,\ldots,c_n\in\Bbb R$ dan
$a_0,\ldots,a_n,b_0,\ldots,b_n\in\BbbR^r$; yakni, $\overline A$
merupakan himpunan tertutup yang memuat
$\{(h\restr X)^{\ssbullet}:h\in S_0\}$, dan merupakan seluruh
$L^p(\mu)$, menurut 244Hb.
}%end of proof of 244I

\leader{244J}{Dualitas dalam ruang $L^p$} Misalkan $(X,\Sigma,\mu)$
sebarang ruang ukur, dan $p\in\ooint{1,\infty}$.
Tetapkan $q=p/(p-1)$\cmmnt{; perhatikan bahwa
${1\over p}+{1\over q}=1$ dan $p=q/(q-1)$; hubungan antara $p$ dan $q$
simetris}. Sekarang $u\times v\in L^1(\mu)$ dan
$\|u\times v\|_1\le\|u\|_p\|v\|_q$ setiap kali
$u\in L^p=L^p(\mu)$ dan $v\in L^q=L^q(\mu)$\cmmnt{ (244E)}.
Akibatnya kita mempunyai suatu operator linear terbatas $T$ dari $L^q$
ke dual ruang bernorma $(L^p)^*$ dari $L^p$, yang diberikan dengan
menuliskan

\Centerline{$(Tv)(u)=\int u\times v$}

\noindent untuk semua $u\in L^p$ dan $v\in L^q$\cmmnt{, persis seperti
dalam 243F}.

\leader{244K}{Teorema} Misalkan $(X,\Sigma,\mu)$ suatu ruang ukur, dan
$p\in\ooint{1,\infty}$; tetapkan $q=p/(p-1)$. Maka peta kanonik
$T:L^q(\mu)\to L^p(\mu)^*$\cmmnt{, yang diuraikan dalam 244J,}
merupakan isomorfisme ruang bernorma.

\cmmnt{\medskip

\noindent{\bf Catatan} Mungkin saya perlu mengingatkan siapa pun yang
membaca bagian ini untuk mempelajari $L^2$ bahwa teori dasar ruang Hilbert
memberikan teorema ini dalam kasus $p=q=2$ tanpa memerlukan argumen yang
lebih umum yang diberikan di bawah (lihat 244N, 244Yk).
}%end of comment

\proof{ Kita mengetahui bahwa $T$ merupakan operator linear terbatas
dengan norma paling besar $1$; saya perlu menunjukkan (i) bahwa $T$
sebenarnya suatu isometri (yakni, bahwa $\|Tv\|=\|v\|_q$ untuk setiap
$v\in L^q$), yang sekaligus akan menunjukkan bahwa $T$ injektif (ii)
bahwa $T$ surjektif, yang merupakan bagian yang benar-benar substansial
dari teorema tersebut.

\medskip

{\bf (a)} Jika $v\in L^q$, maka menurut 244Fa (dengan mengingat bahwa
$p=q/(q-1)$) terdapat $u\in L^p$ sedemikian sehingga $\|u\|_p\le 1$ dan
$\int u\times v=\|v\|_q$; jadi
$\|Tv\|\ge (Tv)(u)=\|v\|_q$. Karena kita telah mengetahui bahwa
$\|Tv\|\le\|v\|_q$, kita mempunyai $\|Tv\|=\|v\|_q$ untuk setiap $v$,
dan $T$ merupakan isometri.

\medskip

{\bf (b)} Karena itu, sisa bukti akan dikhususkan untuk menunjukkan
bahwa $T:L^q\to(L^p)^*$ surjektif. Tetapkan $h\in (L^p)^*$ dengan
$\|h\|=1$.

Saya perlu menunjukkan bahwa $h$ `hidup pada' suatu gabungan terhitung
dari himpunan-himpunan berukuran hingga dalam $X$, dalam arti berikut:
terdapat suatu barisan tak menurun $\sequencen{E_n}$ dari himpunan
berukuran hingga sedemikian sehingga $h(f^{\ssbullet})=0$ setiap kali
$f\in\eusm L^p$ dan $f(x)=0$ untuk
$x\in\bigcup_{n\in\Bbb N}E_n$. \Prf\ Pilih suatu barisan
$\sequencen{u_n}$ dalam $L^p$ sedemikian sehingga $\|u_n\|_p\le 1$ untuk
setiap $n$ dan $\lim_{n\to\infty}h(u_n)=\|h\|=1$.
Untuk setiap $n$, nyatakan $u_n$ sebagai $f_n^{\ssbullet}$, dengan
$f_n:X\to\Bbb R$ suatu fungsi terukur. Tetapkan

\Centerline{$E_n=\{x:\sum_{k=0}^n|f_k(x)|^p\ge 2^{-n}\}$}

\noindent untuk $n\in \Bbb N$; karena $|f_k|^p$ terukur dan terintegralkan
serta berdomain $X$ untuk setiap $k$, $E_n\in \Sigma$ dan
$\mu E_n<\infty$ untuk setiap $n$.

Sekarang andaikan $f\in\eusm L^p(X)$ dan $f(x)=0$ untuk
$x\in\bigcup_{n\in\Bbb N}E_n$; tetapkan
$u=f^{\ssbullet}\in L^p$. \Quer\ Andaikan, jika mungkin, bahwa
$h(u)\ne 0$, dan tinjau $h(cu)$, dengan

\Centerline{$\sgn c=\sgn h(u)$,\quad
$0<|c|<(p\,|h(u)|\,\|u\|_p^{-p})^{1/(p-1)}$.}

\noindent (Tentu saja $\|u\|_p\ne 0$ jika $h(u)\ne 0$.) Untuk setiap
$n$, kita mempunyai

\Centerline{$\{x:f_n(x)\ne 0\}\subseteq\bigcup_{m\in\Bbb
N}E_m\subseteq\{x:f(x)=0\}$,}

\noindent sehingga $|f_n+cf|^p=|f_n|^p+|cf|^p$ dan

\Centerline{$h(u_n+cu)\le\|u_n+cu\|_p=(\|u_n\|_p^p+\|cu\|_p^p)^{1/p}
\le(1+|c|^p\|u\|_p^p)^{1/p}$.}

\noindent Dengan mengambil $n\to\infty$,

\Centerline{$1+ch(u)\le(1+|c|^p\|u\|_p^p)^{1/p}$.}

\noindent
Karena $\sgn c=\sgn h(u)$, $ch(u)=|c||h(u)|$ dan kita mempunyai

\Centerline{$1+p|c||h(u)|\le(1+ch(u))^p\le 1+|c|^p\|u\|_p^p$,}

\noindent sehingga

\Centerline{$p|h(u)|\le |c|^{p-1}\|u\|_p^p<p|h(u)|$}

\noindent menurut pemilihan $c$; ini mustahil.\ \Bang

Ini berarti bahwa $h(f^{\ssbullet})=0$ setiap kali $f:X\to\Bbb R$
terukur, termasuk dalam $\eusm L^p$, dan nol pada
$\bigcup_{n\in\Bbb N}E_n$.\ \Qed

\medskip

{\bf (c)} Tetapkan $H_n=E_n\setminus\bigcup_{k<n}E_k$ untuk setiap
$n\in\Bbb N$; maka $\sequencen{H_n}$ merupakan barisan saling lepas dari
himpunan-himpunan berukuran hingga. Sekarang
$h(u)=\sum_{n=0}^{\infty}h(u\times(\chi H_n)^{\ssbullet})$ untuk setiap
$u\in L^p$. \Prf\ Nyatakan $u$ sebagai $f^{\ssbullet}$, dengan
$f:X\to\Bbb R$ suatu fungsi terukur. Tetapkan
$f_n=f\times\chi H_n$ untuk setiap $n$,
$g=f\times\chi(X\setminus\bigcup_{n\in\Bbb N}H_n)$; maka
$h(g^{\ssbullet})=0$, menurut (a), karena
$\bigcup_{n\in\Bbb N}H_n=\bigcup_{n\in\Bbb N}E_n$. Tinjau

\Centerline{$g_n=g+\sum_{k=0}^nf_k\in \eusm L^p$}

\noindent untuk setiap $n$. Maka $\lim_{n\to\infty}f-g_n=\tbf{0}$, dan

\Centerline{$|f-g_n|^p\le|f|^p\in\eusm L^1$}

\noindent untuk setiap $n$, sehingga menurut Lema Fatou atau Teorema
Konvergensi Terdominasi Lebesgue,

\Centerline{$\lim_{n\to\infty}\int|f-g_n|^p=0$,}

\noindent dan

$$\eqalign{\lim_{n\to\infty}\|u-g^{\ssbullet}-\sum_{k=0}^nu\times(\chi
H_k)^{\ssbullet}\|_p
&=\lim_{n\to\infty}\|u-g_n^{\ssbullet}\|_p\cr
&=\lim_{n\to\infty}\bigl(\int|f-g_n|^p\bigr)^{1/p}
=0,\cr}$$

\noindent yaitu,

\Centerline{$u
=g^{\ssbullet}+\sum_{k=0}^{\infty}u\times\chi H_k^{\ssbullet}$}

\noindent dalam $L^p$. Karena $h:L^p\to\Bbb R$ linear dan kontinu,
diperoleh

\Centerline{$h(u)
=h(g^{\ssbullet})+\sum_{k=0}^{\infty}h(u\times\chi H_k^{\ssbullet})
=\sum_{k=0}^{\infty}h(u\times\chi H_k^{\ssbullet})$,}

\noindent sebagaimana dinyatakan.\ \Qed

\medskip

{\bf (d)} Untuk setiap $n\in\Bbb N$, definisikan
$\nu_n:\Sigma\to\Bbb R$ dengan menetapkan

\Centerline{$\nu_nF=h(\chi(F\cap H_n)^{\ssbullet})$}

\noindent untuk setiap $F\in\Sigma$. (Perhatikan bahwa $\nu_nF$ selalu
terdefinisi karena $\mu(F\cap H_n)<\infty$, sehingga

\Centerline{$\|\chi(F\cap H_n)\|_p=\mu(F\cap H_n)^{1/p}<\infty$.)}

\noindent Maka $\nu_n\emptyset=h(0)=0$, dan jika
$\sequence{k}{F_k}$ suatu barisan saling lepas dalam $\Sigma$,

\Centerline{$\|\chi(\bigcup_{k\in\Bbb N}H_n\cap F_k)
-\sum_{k=0}^m\chi(H_n\cap F_k)\|_p
=\mu(H_n\cap\bigcup_{k=m+1}^{\infty}F_k)^{1/p}\to 0$}

\noindent ketika $m\to\infty$, sehingga

\Centerline{$\nu_n(\bigcup_{k\in\Bbb N}F_k)
=\sum_{k=0}^{\infty}\nu_nF_k$.}

\noindent Jadi $\nu_n$ aditif terhitung. Lebih lanjut,
$|\nu_nF|\le\mu(H_n\cap F)^{1/p}$ untuk setiap $F\in\Sigma$, sehingga
$\nu_n$ benar-benar kontinu dalam arti 232Ab.

Karena itu terdapat suatu fungsi terintegralkan $g_n$ sedemikian sehingga
$\nu_nF=\int_Fg_n$ untuk setiap $F\in\Sigma$; mari kita andaikan bahwa
$g_n$ terukur dan terdefinisi pada seluruh $X$. Tetapkan $g(x)=g_n(x)$
setiap kali $n\in\Bbb N$ dan $x\in H_n$, serta $g(x)=0$ untuk
$x\in X\setminus\bigcup_{n\in\Bbb N}H_n$.

\medskip

{\bf (e)} $g=\sum_{n=0}^{\infty}g_n\times\chi H_n$ terukur dan
mempunyai sifat bahwa $\int_Fg=h(\chi F^{\ssbullet})$ setiap kali
$n\in\Bbb N$ dan $F$ suatu subhimpunan terukur dari $H_n$; akibatnya
$\int_Fg=h(\chi F^{\ssbullet})$ setiap kali $n\in\Bbb N$ dan $F$ suatu
subhimpunan terukur dari $E_n=\bigcup_{k\le n}H_k$. Tetapkan
$G=\{x:g(x)>0\}\subseteq\bigcup_{n\in\Bbb N}E_n$. Jika $F\subseteq G$
dan $\mu F<\infty$, maka

\Centerline{$\lim_{n\to\infty}\int g\times\chi(F\cap E_n)
\le\sup_{n\in\Bbb N}h(\chi(F\cap E_n)^{\ssbullet})
\le\sup_{n\in\Bbb N}\|\chi(F\cap E_n)\|_p
=(\mu F)^{1/p}$,}

\noindent sehingga menurut teorema B.Levi,

\Centerline{$\int_Fg=\int g\times\chi F
=\lim_{n\to\infty}\int g\times\chi(F\cap E_n)$}

\noindent ada. Demikian pula, $\int_Fg$ ada jika
$F\subseteq\{x:g(x)<0\}$ mempunyai ukuran hingga; sedangkan tentu
$\int_Fg$ ada jika $F\subseteq\{x:g(x)=0\}$.
Dengan demikian $\int_Fg$ ada untuk setiap himpunan $F$ berukuran hingga.
Selain itu, menurut Teorema Konvergensi Terdominasi Lebesgue,

\Centerline{$\int_Fg=\lim_{n\to\infty}\int_{F\cap E_n}g
=\lim_{n\to\infty}h(\chi(F\cap E_n)^{\ssbullet})
=\sum_{n=0}^{\infty}h(\chi(F\cap H_n)^{\ssbullet})
=h(\chi F^{\ssbullet})$}

\noindent untuk $F$ demikian, menurut (c) di atas. Segera diperoleh bahwa

\Centerline{$\int g\times f=h(f^{\ssbullet})$}

\noindent untuk setiap fungsi sederhana $f:X\to\Bbb R$.

\medskip

{\bf (f)} Sekarang $g\in L^q$. \Prf\ (i) Kita telah mengetahui bahwa
$|g|^q:X\to\Bbb R$ terukur, karena $g$ terukur dan
$a\mapsto|a|^q$ kontinu. (ii) Andaikan $f$ suatu fungsi sederhana
nonnegatif dan $f\leae|g|^q$. Maka $f^{1/p}$ suatu fungsi sederhana, dan
$\sgn g$ terukur serta hanya mengambil nilai $0$, $1$ dan $-1$, sehingga
$f_1=f^{1/p}\times\sgn g$ sederhana. Kita melihat bahwa
$\int|f_1|^p=\int f$, sehingga $\|f_1\|_p=(\int f)^{1/p}$. Dengan
demikian

$$\eqalignno{(\int f)^{1/p}&\ge h(f_1^{\ssbullet})
=\int g\times f_1=\int|g\times f^{1/p}|\cr
&\ge\int f^{1/q}\times f^{1/p}\cr
\noalign{\noindent (karena $0\le f^{1/q}\leae|g|$)}
&=\int f,\cr}$$

\noindent dan kita harus mempunyai $\int f\le 1$. (iii) Jadi

\Centerline{$\sup\{\int f: f$ suatu fungsi sederhana nonnegatif,
$f\leae |g|^q\}\le 1<\infty$.}

\noindent Namun sekarang perhatikan bahwa jika $\epsilon>0$, maka

\Centerline{$\{x:|g(x)|^q\ge\epsilon\}
=\bigcup_{n\in\Bbb N}\{x:x\in E_n,\,|g(x)|^q\ge\epsilon\}$,}

\noindent dan untuk setiap $n\in\Bbb N$,

\Centerline{$\mu\{x:x\in E_n,\,|g(x)|^q\ge\epsilon\}
\le{1\over{\epsilon}}$,}

\noindent karena
$f=\epsilon\chi\{x:x\in E_n,\,|g(x)|^q\ge\epsilon\}$ merupakan fungsi
sederhana yang kurang dari atau sama dengan $|g|^q$, sehingga integralnya
paling besar $1$. Dengan demikian

\Centerline{$\mu\{x:|g(x)|^q\ge\epsilon\}
=\sup_{n\in\Bbb N}\mu\{x:x\in E_n,\,|g(x)|^q\ge\epsilon\}
\le{1\over{\epsilon}}<\infty$.}

\noindent Jadi $|g|^q$ terintegralkan, menurut kriteria dalam 122Ja.\ \Qed

\medskip

{\bf (g)} Karena itu kita dapat berbicara tentang
$h_1=T(g^{\ssbullet})\in (L^p)^*$, dan kita mengetahui bahwa ia sama
dengan $h$ pada anggota-anggota $L^p$ berbentuk $f^{\ssbullet}$, dengan
$f$ suatu fungsi sederhana. Namun anggota-anggota ini membentuk
subhimpunan padat dari $L^p$, menurut 244Ha, dan baik $h$ maupun $h_1$
kontinu, sehingga $h=h_1$, menurut 2A3Uc, dan $h$ merupakan suatu nilai
$T$. Argumen sejauh ini mengandaikan bahwa $\|h\|=1$. Namun tentu saja
setiap anggota tak nol dari $(L^p)^*$ merupakan kelipatan skalar dari
suatu unsur bernorma $1$, sehingga merupakan suatu nilai $T$. Jadi
$T:L^q\to(L^p)^*$ memang surjektif, dan karena itu merupakan isomorfisme
isometrik, sebagaimana dinyatakan.
}%end of proof of 244K

\vleader{48pt}{244L}{}\cmmnt{ Melanjutkan topik-topik yang sama seperti
dalam \S\S242 dan 243, saya beralih ke kelengkapan urutan $L^p$.

\medskip

\noindent}{\bf Teorema} Misalkan $(X,\Sigma,\mu)$ sebarang ruang ukur,
dan $p\in\coint{1,\infty}$. Maka $L^p=L^p(\mu)$ lengkap Dedekind.

\proof{ Saya menggunakan 242H. Misalkan $A\subseteq L^p$ suatu himpunan
tak kosong yang terbatas di atas dalam $L^p$. Tetapkan $u_0\in A$ dan
tetapkan

\Centerline{$A'=\{u_0\vee u:u\in A\}$,}

\noindent sehingga $A'$ mempunyai batas-batas atas yang sama dengan $A$
dan terbatas di bawah oleh $u_0$. Dengan menetapkan suatu batas atas
$w_0$ dari $A$ dalam $L^p$, maka $u_0\le u\le w_0$ untuk setiap
$u\in A'$. Tetapkan

\Centerline{$B=\{(u-u_0)^p:u\in A'\}$.}

\noindent Maka

\Centerline{$0\le v\le(w_0-u_0)^p\in L^1=L^1(\mu)$}

\noindent untuk setiap $v\in B$, sehingga $B$ merupakan subhimpunan tak
kosong dari $L^1$ yang terbatas di atas dalam $L^1$, dan karena itu
mempunyai batas atas terkecil $v_1$ dalam $L^1$. Sekarang
$v_1^{1/p}\in L^p$; tinjau $w_1=u_0+v_1^{1/p}$.
Jika $u\in A'$, maka $(u-u_0)^p\le v_1$, sehingga
$u-u_0\le v_1^{1/p}$ dan $u\le w_1$; jadi $w_1$ suatu batas atas untuk
$A'$. Jika $w\in L^p$ suatu batas atas untuk $A'$, maka
$u-u_0\le w-u_0$ dan $(u-u_0)^p\le(w-u_0)^p$ untuk setiap $u\in A'$,
sehingga $(w-u_0)^p$ suatu batas atas untuk $B$ dan
$v_1\le(w-u_0)^p$, $v_1^{1/p}\le w-u_0$, serta $w_1\le w$. Jadi
$w_1=\sup A'=\sup A$ dalam $L^p$. Karena $A$ sembarang, $L^p$ lengkap
Dedekind.
}%end of proof of 244L

\leader{244M}{}\cmmnt{ Seperti dalam dua bagian terakhir, teori
ekspektasi bersyarat patut ditinjau kembali.

\medskip

\noindent}{\bf Teorema} Misalkan $(X,\Sigma,\mu)$ suatu ruang
probabilitas, dan $\Tau$ suatu subaljabar-$\sigma$ dari $\Sigma$.
Ambil $p\in[1,\infty]$. Pandang $L^0(\mu\restrp\Tau)$ sebagai subruang
dari $L^0=L^0(\mu)$\cmmnt{, seperti dalam 242Jh}, sehingga
$L^p(\mu\restrp\Tau)$ menjadi
$L^p(\mu)\cap L^0(\mu\restrp\Tau)$. Misalkan
$P:L^1(\mu)\to L^1(\mu\restrp\Tau)$ operator ekspektasi
bersyarat\cmmnt{, seperti diuraikan dalam 242Jd}. Maka setiap kali
$u\in L^p=L^p(\mu)$, $|Pu|^p\le P(|u|^p)$, sehingga
$Pu\in L^p(\mu\restrp\Tau)$ dan $\|Pu\|_p\le\|u\|_p$. Selain itu,
$P[L^p]=L^p(\mu\restrp\Tau)$.

\proof{ Untuk $p=\infty$, ini adalah 243Jb, jadi selanjutnya saya
andaikan bahwa $p<\infty$. Mengenai identifikasi
$L^p(\mu\restrp\Tau)$ dengan $L^p\cap L^0(\mu\restrp\Tau)$, jika
$S:L^0(\mu\restrp\Tau)\to L^0$ merupakan pembenaman kanonik yang
diuraikan dalam 242J, kita mempunyai $|Su|^p=S(|u|^p)$ untuk setiap
$u\in L^0(\mu\restrp\Tau)$, sehingga $Su\in L^p$ jika dan hanya jika
$|u|^p\in L^1(\mu\restrp\Tau)$ jika dan hanya jika
$u\in L^p(\mu\restrp\Tau)$.

Tetapkan $\phi(t)=|t|^p$ untuk $t\in\Bbb R$; maka $\phi$ suatu fungsi
konveks (karena kontinu mutlak pada sebarang interval terbatas, dan
turunannya tak menurun), dan $|u|^p=\bar\phi(u)$ untuk setiap
$u\in L^0=L^0(\mu)$, dengan $\bar\phi$ didefinisikan seperti dalam 241I.
Sekarang jika $u\in L^p=L^p(\mu)$, tentu kita mempunyai $u\in L^1$
(karena $|u|\le|u|^p\vee(\chi X)^{\ssbullet}$, atau dengan cara lain);
jadi 242K memberi tahu kita bahwa $|Pu|^p\le P|u|^p$. Namun ini berarti
bahwa $Pu\in L^p\cap L^1(\mu\restrp\Tau)=L^p(\mu\restrp\Tau)$, dan

\Centerline{$\|Pu\|_p=(\int|Pu|^p)^{1/p}\le(\int P|u|^p)^{1/p}
=(\int|u|^p)^{1/p}=\|u\|_p$,}

\noindent sebagaimana dinyatakan. Jika $u\in L^p(\mu\restrp\Tau)$,
maka $Pu=u$, sehingga $P[L^p]$ adalah seluruh
$L^p(\mu\restrp\Tau)$.
}%end of proof of 244M


\leader{244N}{Ruang $L^2$ (a)}\cmmnt{ Seperti telah saya nyatakan,
ruang fungsi yang benar-benar penting ialah $L^0$, $L^1$, $L^2$ dan
$L^{\infty}$.} $L^2$ mempunyai sifat khusus sebagai ruang hasil kali
dalam; jika $(X,\Sigma,\mu)$ sebarang ruang ukur dan
$u$, $v\in L^2=L^2(\mu)$, maka $u\times v\in L^1(\mu)$\cmmnt{, menurut
244Eb}, dan kita dapat menuliskan $\innerprod{u}{v}=\int u\times v$.
Ini menjadikan $L^2$ suatu ruang hasil kali dalam real\prooflet{ (karena

\Centerline{$\innerprod{u_1+u_2}{v}
=\innerprod{u_1}{v}+\innerprod{u_2}{v}$,
\quad
$\innerprod{cu}{v}=c\innerprod{u}{v}$,
\quad
$\innerprod{u}{v}=\innerprod{v}{u}$,}

\Centerline{$\innerprod{u}{u}\ge 0$,
\quad $u=0$ setiap kali $\innerprod{u}{u}=0$}

\noindent untuk semua $u$, $u_1$, $u_2$, $v\in L^2$ dan $c\in\Bbb R$)}
dan normanya $\|\,\|_2$ merupakan norma yang terkait\cmmnt{ (karena
$\|u\|_2=\sqrt{\innerprod{u}{u}}$ setiap kali $u\in L^2$)}.
\cmmnt{Karena} $L^2$\cmmnt{ lengkap (244G), ruang ini} merupakan ruang
Hilbert real. \cmmnt{Fakta bahwa ruang ini dapat diidentifikasi dengan
dualnya sendiri (244K) tentu dapat diturunkan dari sini.}

Saya akan menggunakan frasa `{\bf terintegralkan-kuadrat}' untuk
menyebut fungsi-fungsi dalam $\eusm L^2(\mu)$.

\spheader 244Nb \cmmnt{Ekspektasi bersyarat mengambil bentuk khusus dalam
kasus $L^2$.} Misalkan $(X,\Sigma,\mu)$ suatu ruang probabilitas,
$\Tau$ suatu subaljabar-$\sigma$ dari $\Sigma$, dan
$P:L^1=L^1(\mu)\to L^1(\mu\restrp\Tau)\subseteq L^1$ operator ekspektasi
bersyarat yang sesuai. Maka $P[L^2]\subseteq L^2$, dengan
$L^2=L^2(\mu)$\cmmnt{ (244M)}, sehingga kita mempunyai suatu operator
$P_2=P\restr L^2$ dari $L^2$ ke dirinya sendiri. Sekarang $P_2$ merupakan
proyeksi ortogonal dan kernelnya ialah
$\{u:u\in L^2,\,\int_Fu=0$ untuk setiap $F\in\Tau\}$.
\prooflet{\Prf\ (i) Jika $u\in L^1$, maka $Pu=0$ jika dan hanya jika
$\int_Fu=0$ untuk setiap $F\in\Tau$ (bandingkan 242Je); jadi tentu kernel
$P_2$ adalah himpunan yang diuraikan. (ii) Karena $P^2=P$, $P_2$ juga
merupakan suatu proyeksi; karena $P_2$ mempunyai norma paling besar $1$
(244M), dan karena itu kontinu,

\Centerline{$U=P_2[L^2]=L^2(\mu\restrp\Tau)
=\{u:u\in L^2,\,P_2u=u\}$,\quad$V=\{u:P_2u=0\}$}

\noindent merupakan subruang-subruang linear tertutup dari $L^2$
sedemikian sehingga $U\oplus V=L^2$.
(iii) Sekarang andaikan $u\in U$ dan $v\in V$. Maka $P|v|\in L^2$,
sehingga $u\times P|v|\in L^1$ dan $P(u\times v)=u\times Pv$, menurut
242L. Dengan demikian

\Centerline{$\innerprod{u}{v}=\int u\times v
=\int P(u\times v)=\int u\times Pv=0$.}

\noindent Jadi $U$ dan $V$ merupakan subruang ortogonal dari $L^2$,
yang kita maksud dengan mengatakan bahwa $P_2$ merupakan proyeksi
ortogonal. (Sebagian pembaca akan mengetahui bahwa setiap proyeksi
bernorma paling besar $1$ pada ruang hasil kali dalam bersifat
ortogonal.)\ \Qed
}%end of prooflet

\leader{*244O}{}\dvAnew{2009}\cmmnt{ Ini bukan tempat untuk pembahasan
terperinci tentang geometri ruang-ruang $L^p$. Namun terdapat suatu fakta
yang sangat penting mengenai bentuk bola satuan yang dapat dijangkau
dengan metode-metode yang tersedia bagi kita di sini.

\medskip

\noindent}{\bf Teorema}\cmmnt{ ({\smc Clarkson 36})}
Andaikan $p\in\ooint{1,\infty}$ dan $(X,\Sigma,\mu)$ suatu ruang ukur.
Maka $L^p=L^p(\mu)$ konveks seragam\cmmnt{ (definisi: 2A4K)}.

\proof{({\smc Hanner 56}, {\smc Naor 04})

\medskip

{\bf (a)(i)} Untuk $0<t\le 1$ dan $a$, $b\in\Bbb R$, tetapkan

\Centerline{$\phi_0(t)=(1+t)^{p-1}+(1-t)^{p-1}$,}

\Centerline{$\phi_1(t)=\Bover{(1+t)^{p-1}-(1-t)^{p-1}}{t^{p-1}}
=(\Bover1t+1)^{p-1}-(\Bover1t-1)^{p-1}$,}

\Centerline{$\psi_{ab}(t)=|a|^p\phi_0(t)+|b|^p\phi_1(t)$,}

\Centerline{$\phi_2(b)=(1+b)^p+|1-b|^p$.}

\medskip

\quad{\bf (ii)} Kita mempunyai

$$\eqalignno{\phi_0'(t)
&=(p-1)((1+t)^{p-2}-(1-t)^{p-2}),\text{ yang bertanda sama dengan }p-2,\cr
\noalign{\noindent (tentu saja nilainya nol jika $p=2$),}
\phi_1'(t)
&=-\Bover{p-1}{t^2}((\Bover1t+1)^{p-2}-(\Bover1t-1)^{p-2})\cr
&=-\Bover{p-1}{t^p}((1+t)^{p-2}-(1-t)^{p-2})
=-\Bover1{t^p}\phi_0'(t)\cr}$$

\noindent untuk setiap $t\in\ooint{0,1}$. Dengan demikian
$\phi_0'-\phi_1'$ mempunyai tanda yang sama dengan $p-2$ di mana-mana
pada $\ooint{0,1}$. Selain itu,

\Centerline{$\phi_0(1)=2^{p-1}=\phi_1(1)$,}

\noindent sehingga $\phi_0-\phi_1$ mempunyai tanda yang sama dengan
$2-p$ di mana-mana pada $\ocint{0,1}$.

\medskip

\quad{\bf (iii)} $\phi_2$ naik secara ketat pada
$\coint{0,\infty}$. \Prf\ Untuk $b>0$,

$$\eqalign{\phi_2'(b)
&=p((1+b)^{p-1}-(1-b)^{p-1})>0\text{ jika }b\le 1,\cr
&=p((1+b)^{p-1}+(b-1)^{p-1})>0\text{ jika }b\ge 1.
\text{ \Qed}\cr}$$

\medskip

\quad{\bf (iv)} Jika $0<b\le a$, maka

$$\eqalignno{\psi_{ab}(\Bover{b}{a})
&=a^p\phi_0(\Bover{b}{a})+b^p\phi_1(\Bover{b}{a})\cr
&=a^p(1+\Bover{b}{a})^{p-1}+a^p(1-\Bover{b}{a})^{p-1}
   +b^p(\Bover{a}{b}+1)^{p-1}-b^p(\Bover{a}{b}-1)^{p-1}\cr
&=a(a+b)^{p-1}+a(a-b)^{p-1}+b(a+b)^{p-1}-b(a-b)^{p-1}\cr
&=(a+b)^p+(a-b)^p
=(a+b)^p+|a-b|^p.&(\dagger)\cr}$$

\noindent Selain itu,
$\psi_{ab}'(t)=(a^p-\Bover{b^p}{t^p})\phi_0'(t)$ bertanda sama dengan
$2-p$ jika $0<t<\Bover{b}{a}$ dan bertanda sama dengan $p-2$ jika
$\Bover{b}{a}<t<1$. Dengan demikian

\inset{----- jika $1<p\le 2$,
$\psi_{ab}(t)\le\psi_{ab}(\Bover{b}{a})=(a+b)^p+|a-b|^p$ untuk setiap
$t\in\ocint{0,1}$,

----- jika $p\ge 2$,
$\psi_{ab}(t)\ge\psi_{ab}(\Bover{b}{a})=(a+b)^p+|a-b|^p$ untuk setiap
$t\in\ocint{0,1}$.}

\medskip

\quad{\bf (v)} Sekarang tinjau kasus $0<a\le b$. Jika $1<p\le 2$,

$$\eqalignno{\psi_{ab}(t)
&=a^p\phi_0(t)+b^p\phi_1(t)
\le a^p\phi_0(t)+b^p\phi_1(t)+(b^p-a^p)(\phi_0(t)-\phi_1(t))\cr
\displaycause{menurut (ii)}
&=b^p\phi_0(t)+a^p\phi_1(t)
\le(b+a)^p+(b-a)^p
=(a+b)^p+|a-b|^p\cr}$$

\noindent untuk setiap $t\in\ocint{0,1}$. Sebaliknya, jika $p\ge 2$,

$$\eqalignno{\psi_{ab}(t)
&=a^p\phi_0(t)+b^p\phi_1(t)
\ge a^p\phi_0(t)+b^p\phi_1(t)+(b^p-a^p)(\phi_0(t)-\phi_1(t))\cr
&=b^p\phi_0(t)+a^p\phi_1(t)
\ge(a+b)^p+|a-b|^p\cr}$$

\noindent untuk setiap $t$.

\medskip

\quad{\bf (vi)} Jadi kita mempunyai pertidaksamaan

$$\eqalign{\psi_{ab}(t)
&\le|a+b|^p+|a-b|^p\text{ jika }p\in\ocint{1,2},\cr
&\ge|a+b|^p+|a-b|^p\text{ jika }p\in\coint{2,\infty}\cr}\eqno{(*)}$$

\noindent setiap kali $t\in\ocint{0,1}$ dan
$a$, $b\in\ooint{0,\infty}$. Karena
$(a,b)\mapsto\psi_{ab}(t)$ kontinu untuk setiap $t$, pertidaksamaan yang
sama berlaku untuk semua $a$, $b\in\coint{0,\infty}$. Dan karena

\Centerline{$\psi_{ab}(t)=\psi_{|a|,|b|}(t)$,
\quad$|a+b|^p+|a-b|^p=\big||a|+|b|\bigr|^p+\bigl||a|-|b|\bigr|^p$}

\noindent untuk semua $a$, $b\in\Bbb R$ dan $t\in\ocint{0,1}$,
pertidaksamaan (*) berlaku untuk semua $a$, $b\in\Bbb R$ dan
$t\in\ocint{0,1}$.

\medskip

{\bf (b)} Andaikan $p\ge 2$.

\medskip

\quad{\bf (i)}

\Centerline{$\|u+v\|_p^p+\|u-v\|_p^p
\le(\|u\|_p+\|v\|_p)^p+|\|u\|_p-\|v\|_p|^p$}

\noindent untuk semua $u$, $v\in L^p$. \Prf\ Pertama tinjau kasus
$0<\|v\|_p\le\|u\|_p$. Misalkan $f$, $g:X\to\Bbb R$ fungsi-fungsi yang
terukur terhadap $\Sigma$ sedemikian sehingga $f^{\ssbullet}=u$ dan
$g^{\ssbullet}=v$. Maka untuk sebarang $t\in\ocint{0,1}$,

$$\eqalignno{\|u+v\|_p^p+\|u-v\|_p^p
&=\int|f(x)+g(x)|^p+|f(x)-g(x)|^p\mu(dx)\cr
&\le\int\psi_{f(x),g(x)}(t)\mu(dx)\cr
\displaycause{menurut pertidaksamaan kedua dalam (*)}
&=\int|f(x)|^p\phi_0(t)+|g(x)|^p\phi_1(t)\mu(dx)
=\|u\|_p^p\phi_0(t)+\|v\|_p^p\phi_1(t).\cr}$$

\noindent Secara khusus, dengan mengambil
$t=\|v\|_p/\|u\|_p$, dan menerapkan ($\dagger$) dari (a-iv),

\Centerline{$\|u+v\|_p^p+\|u-v\|_p^p
\le(\|u\|_p+\|v\|_p)^p+|\|u\|_p-\|v\|_p|^p$.}

\noindent Tentu saja hasil tersebut juga benar jika
$0<\|u\|_p\le\|v\|_p$, dan kasus ketika $u$ atau $v$ bernilai nol
bersifat sepele.\ \Qed

\medskip

\quad{\bf (ii)} Misalkan $\epsilon\in\ocint{0,2}$. Tetapkan
$\delta=2-(2^p-\epsilon^p)^{1/p}>0$. Jika $u$, $v\in L^p$,
$\|u\|_p=\|v\|_p=1$ dan $\|u-v\|_p\ge\epsilon$, maka

\Centerline{$\|u+v\|_p^p+\epsilon^p
\le\|u+v\|_p^p+\|u-v\|_p^p
\le(\|u\|_p+\|v\|_p)^p+|\|u\|_p-\|v\|_p|^p
=2^p$,}

\noindent sehingga $\|u+v\|_p\le(2^p-\epsilon^p)^{1/p}=2-\delta$.
Karena $u$, $v$ dan $\epsilon$ sembarang, $L^p$ konveks seragam.

\wheader{*244O}{6}{2}{2}{36pt}

{\bf (c)} Selanjutnya andaikan $p\in\ocint{1,2}$.

\medskip

\quad{\bf (i)}

\Centerline{$(\|u\|_p+\|v\|_p)^p+|\|u\|_p-\|v\|_p|^p
\le\|u+v\|_p^p+\|u-v\|_p^p$}

\noindent untuk semua $u$, $v\in L^p$. \Prf\ Kita dapat mengulangi
semua gagasan, dan sebagian besar rumus, dari (b-i). Seperti sebelumnya,
mulailah dengan kasus $0<\|v\|_p\le\|u\|_p$. Misalkan
$f$, $g:X\to\Bbb R$ fungsi-fungsi yang terukur terhadap $\Sigma$
sedemikian sehingga $f^{\ssbullet}=u$ dan $g^{\ssbullet}=v$.
Dengan mengambil $t=\|v\|_p/\|u\|_p$,

$$\eqalignno{\|u+v\|_p^p+\|u-v\|_p^p
&=\int|f(x)+g(x)|^p+|f(x)-g(x)|^p\mu(dx)\cr
&\ge\int\psi_{f(x),g(x)}(t)\mu(dx)\cr
\displaycause{menurut pertidaksamaan pertama dalam (*)}
&=\|u\|_p^p\phi_0(t)+\|v\|_p^p\phi_1(t)
=(\|u\|_p+\|v\|_p)^p+|\|u\|_p-\|v\|_p|^p.\cr}$$

\noindent Demikian pula jika $0<\|u\|_p\le\|v\|_p$, dan kasus ketika
$u$ atau $v$ bernilai nol bersifat sepele.\ \Qed

\medskip

\quad{\bf (ii)} Misalkan $\epsilon>0$. Tetapkan
$\gamma=\phi_2(\bover{\epsilon}2)>2$ (lihat (a-iii) di atas) dan
$\delta=2\bigl(1-(\Bover{2}{\gamma})^{1/p}\bigr)>0$.
Sekarang andaikan bahwa $\|u\|_p=\|v\|_p=1$ dan
$\|u-v\|_p\ge\epsilon$. Maka $\|u+v\|_p\le 2-\delta$. \Prf\ Jika
$u+v=0$, pernyataan ini sepele. Jika tidak, tetapkan
$a=\|u+v\|_p$ dan $b=\|u-v\|_p$. Maka $a\le 2$ dan $b\ge\epsilon$,
sehingga

$$\eqalignno{a^p\gamma
&=a^p\phi_2(\Bover{\epsilon}2)
\le a^p\phi_2(\Bover{b}{a})\cr
\displaycause{sekali lagi menurut (a-iii)}
&=(a+b)^p+|a-b|^p
=(\|u+v\|_p+\|u-v\|_p)^p+|\|u+v\|_p-\|u-v\|_p|^p\cr
&\le\|2u\|_p^p+\|2v\|_p^p\cr
\displaycause{menurut (i) di sini}
&=2^{p+1}\cr}$$

\noindent dan
$a\le 2\bigl(\Bover{2}{\gamma}\bigr)^{1/p}=2-\delta$.\ \QeD\ Karena
$u$, $v$ dan $\epsilon$ sembarang, $L^p$ konveks seragam.
}%end of proof of 244O?

\cmmnt{\medskip

\noindent{\bf Catatan} Pertidaksamaan dalam (b-i) dan (c-i) dari bukti
tersebut disebut {\bf pertidaksamaan Hanner}.}

\leader{244P}{$L^p$ kompleks}\dvAformerly{2{}44O}
Misalkan $(X,\Sigma,\mu)$ sebarang ruang ukur.

\header{244Pa}{\bf (a)} Untuk sebarang $p\in\ooint{1,\infty}$, tetapkan

\Centerline{$\eusm L^p_{\Bbb C}=\eusm L^p_{\Bbb C}(\mu)
=\{f:f\in\eusm L^0_{\Bbb C}(\mu),\,|f|^p$ terintegralkan$\}$,}

$$\eqalign{L^p_{\Bbb C}=L^p_{\Bbb C}(\mu)
&=\{f^{\ssbullet}:f\in\eusm L^p_{\Bbb C}\}\cr
&=\{u:u\in L^0_{\Bbb C}(\mu),\,\Real(u)\in L^p(\mu)\text{ dan }
\Imag(u)\in L^p(\mu)\}\cr
&=\{u:u\in L^0_{\Bbb C}(\mu),\,|u|\in L^p(\mu)\}.\cr}$$

\noindent Maka $L^p_{\Bbb C}$ merupakan subruang linear dari
$L^0_{\Bbb C}(\mu)$. Tetapkan
$\|u\|_p=\||u|\|_p=(\int|u|^p)^{1/p}$ untuk $u\in L^p_{\Bbb C}$.

\header{244Pb}{\bf (b)} Bukti 244E(b-i) berlaku tanpa perubahan untuk
fungsi bernilai kompleks, sehingga dengan mengambil $q=p/(p-1)$ kita
memperoleh

\Centerline{$\|u\times v\|_1\le\|u\|_p\|v\|_q$}

\noindent untuk semua $u\in L^p_{\Bbb C}$, $v\in L^q_{\Bbb C}$. 244Fa
menjadi

\inset{untuk setiap $u\in L^p_{\Bbb C}$ terdapat
$v\in L^q_{\Bbb C}$ sedemikian sehingga $\|v\|_q\le 1$ dan

\Centerline{$\int u\times v=|\int u\times v|=\|u\|_p$\dvro{.}{;}}}

\noindent\cmmnt{bukti yang sama berlaku, jika Anda menghilangkan semua
penyebutan fungsional $\tau$, dan mengizinkan saya menuliskan
$\sgn a=|a|/a$ untuk semua bilangan kompleks tak nol -- mungkin lebih
wajar menuliskan $\overline{\sgn}a$ sebagai pengganti $\sgn a$. Jadi,
seperti sebelumnya, kita mendapati bahwa} $\|\,\|_p$ merupakan suatu
norma. \cmmnt{Kita dapat menggunakan argumen 244G untuk menunjukkan bahwa}
$L^p_{\Bbb C}$ lengkap. \cmmnt{(Sebagai alternatif, perhatikan bahwa
suatu barisan $\sequencen{u_n}$ dalam $L^0_{\Bbb C}$ bersifat Cauchy,
atau konvergen, jika dan hanya jika bagian real dan imajinernya demikian.)}
Ruang $S_{\Bbb C}$ dari kelas-kelas ekuivalensi `fungsi sederhana
bernilai kompleks' padat dalam $L^p_{\Bbb C}$. Jika $X$ subhimpunan dari
$\BbbR^r$ dan $\mu$ ukuran Lebesgue pada $X$, maka ruang kelas-kelas
ekuivalensi fungsi bernilai kompleks kontinu pada $X$ dengan dukungan
terbatas padat dalam $L^p_{\Bbb C}$.

\header{244Pc}{\bf (c)} Peta kanonik
$T:L^q_{\Bbb C}\to (L^p_{\Bbb C})^*$, yang didefinisikan dengan
menuliskan $(Tv)(u)=\int u\times v$, bersifat surjektif\cmmnt{ karena
$T\restr L^q:L^q\to (L^p)^*$ surjektif;} dan\cmmnt{ peta itu} merupakan
isometri\cmmnt{ menurut catatan dalam (b) tepat di atas}. Jadi kita masih
dapat mengidentifikasi $L^q_{\Bbb C}$ dengan $(L^p_{\Bbb C})^*$.

\header{244Pd}{\bf (d)}\cmmnt{ Ketika kita sampai pada bentuk kompleks
pertidaksamaan Jensen, tampaknya diperlukan suatu gagasan baru. Saya
menempatkan ini dalam 242Yk-242Yl. Namun, untuk bentuk kompleks 244M,
argumen yang lebih sederhana sudah memadai.} Jika $(X,\Sigma,\mu)$ suatu
ruang probabilitas, $\Tau$ suatu subaljabar-$\sigma$ dari $\Sigma$ dan
$P:L^1_{\Bbb C}(\mu)\to L^1_{\Bbb C}(\mu\restrp\Tau)$ operator
ekspektasi bersyarat yang sesuai, maka untuk sebarang
$u\in L^p_{\Bbb C}$\cmmnt{ kita akan mempunyai

\Centerline{$|Pu|^p\le(P|u|)^p\le P(|u|^p)$,}

\noindent dengan menerapkan 242Pc dan 244M.
Jadi} $\|Pu\|_p\le\|u\|_p$\cmmnt{, seperti sebelumnya}.

\header{244Pe}{\bf (e)}\cmmnt{ Terdapat satu hal khusus yang muncul pada
$L^2_{\Bbb C}$.} Sekarang kita harus mendefinisikan

\Centerline{$\innerprod{u}{v}=\int u\times\bar v$}

\noindent untuk $u$, $v\in L^2_{\Bbb C}$\cmmnt{, sehingga
$\innerprod{u}{u}=\int|u|^2=\|u\|_2^2$ untuk setiap $u$}; \cmmnt{ini
berarti bahwa} $\innerprod{v}{u}$ merupakan konjugat kompleks dari
$\innerprod{u}{v}$.

\exercises{
\leader{244X}{Latihan dasar $\pmb{>}$(a)}
%\spheader 244Xa
Misalkan $(X,\Sigma,\mu)$ suatu ruang ukur, dan
$(X,\hat\Sigma,\hat\mu)$ pelengkapannya. Tunjukkan bahwa
$\eusm L^p(\hat\mu)=\eusm L^p(\mu)$ dan
$L^p(\hat\mu)=L^p(\mu)$ untuk setiap $p\in[1,\infty]$.
%244A used in 465F

\spheader 244Xb Misalkan $(X,\Sigma,\mu)$ suatu ruang ukur, dan
$1\le p\le q\le r\le\infty$. Tunjukkan bahwa
$L^p(\mu)\cap L^r(\mu)\subseteq L^q(\mu)
\subseteq L^p(\mu)+L^r(\mu)\subseteq L^0(\mu)$. (Lihat pula 244Yh.)
%244A

\spheader 244Xc Misalkan $(X,\Sigma,\mu)$ suatu ruang ukur. Andaikan
$p$, $q$, $r\in[1,\infty]$ dan $\bover1p+\bover1q=\bover1r$, dengan
menetapkan $\bover1{\infty}=0$ seperti biasa. Tunjukkan bahwa
$u\times v\in L^r(\mu)$ dan
$\|u\times v\|_r\le\|u\|_p\|v\|_q$ setiap kali $u\in L^p(\mu)$ dan
$v\in L^q(\mu)$. \Hint{jika $r<\infty$, terapkan pertidaksamaan
H\"older pada $|u|^r\in L^{p/r}$, $|v|^r\in L^{q/r}$.}
%244E

\sqheader 244Xd(i) Misalkan $(X,\Sigma,\mu)$ suatu ruang probabilitas.
Tunjukkan bahwa jika $1\le p\le r\le\infty$, maka
$\|f\|_p\le\|f\|_r$ untuk setiap $f\in{\eusm L}^r(\mu)$.
\Hint{gunakan pertidaksamaan H\"older untuk menunjukkan bahwa
$\int|f|^p\le\||f|^p\|_{r/p}$.} Secara khusus,
${\eusm L}^p(\mu)\supseteq {\eusm L}^r(\mu)$. (ii) Misalkan
$(X,\Sigma,\mu)$ suatu ruang ukur sedemikian sehingga $\mu E\ge 1$
setiap kali $E\in\Sigma$ dan $\mu E>0$. (Ini terjadi, misalnya, ketika
$\mu$ merupakan `ukuran pencacahan' pada $X$.) Tunjukkan bahwa jika
$1\le p\le r\le\infty$, maka $L^p(\mu)\subseteq L^r(\mu)$ dan
$\|u\|_p\ge\|u\|_r$ untuk setiap $u\in L^p(\mu)$.
\Hint{tinjau terlebih dahulu kasus $\|u\|_p=1$.}
%244F

\sqheader 244Xe\dvAformerly{2{}44Xf}
Misalkan $(X,\Sigma,\mu)$ suatu ruang ukur semihingga, dan
$p$, $q\in[1,\infty]$ sedemikian sehingga
$\bover1p+\bover1q=1$. Tunjukkan bahwa jika
$u\in L^0(\mu)\setminus L^p(\mu)$, maka terdapat $v\in L^q(\mu)$
sedemikian sehingga $u\times v\notin L^1(\mu)$.
\Hint{reduksikan ke kasus $u\ge 0$. Tunjukkan bahwa dalam kasus ini,
untuk setiap $n\in\Bbb N$ terdapat $u_n\le u$ sedemikian sehingga
$4^n\le\|u_n\|_p<\infty$; ambil $v_n\in L^q$ sedemikian sehingga
$\|v_n\|_q\le 2^{-n}$ dan $\int u_n\times v_n\ge 2^n$, lalu tetapkan
$v=\sum_{n=0}^{\infty}v_n$.}
%244F

\spheader 244Xf Misalkan
$\langle(X_i,\Sigma_i,\mu_i)\rangle_{i\in I}$ suatu keluarga ruang ukur,
dan $(X,\Sigma,\mu)$ jumlah langsungnya (214L). Ambil sebarang
$p\in\coint{1,\infty}$. Tunjukkan bahwa isomorfisme kanonik antara
$L^0(\mu)$ dan $\prod_{i\in I}L^0(\mu_i)$ (241Xd) menginduksi suatu
isomorfisme antara $L^p(\mu)$ dan subruang

\Centerline{$\{u:u\in\prod_{i\in I}L^p(\mu_i),\,
\|u\|=\bigl(\sum_{i\in I}\|u(i)\|_p^p)^{1/p}<\infty\}$}

\noindent dari $\prod_{i\in I}L^p(\mu_i)$.
%244F

\spheader 244Xg Misalkan $(X,\Sigma,\mu)$ suatu ruang ukur. Tetapkan
$M^{\infty,1}=L^1(\mu)\cap L^{\infty}(\mu)$. Tunjukkan bahwa untuk
$u\in M^{\infty,1}$, fungsi
$p\mapsto\|u\|_p:\coint{1,\infty}\to\coint{0,\infty}$ kontinu, dan bahwa
$\|u\|_{\infty}=\lim_{p\to\infty}\|u\|_p$. \Hint{tinjau terlebih dahulu
kasus ketika $u$ merupakan kelas ekuivalensi dari suatu fungsi sederhana.}
%244F

\spheader 244Xh Misalkan $\mu$ ukuran pencacahan pada $X=\{1,2\}$,
sehingga $\eusm L^0(\mu)=\BbbR^2$ dan $L^p(\mu)=L^0(\mu)$ dapat
diidentifikasi dengan $\BbbR^2$ untuk setiap $p\in[1,\infty]$. Buat
sketsa bola satuan $\{u:\|u\|_p\le 1\}$ dalam $\BbbR^2$ untuk
$p=1$, $\bover32$, $2$, $3$ dan $\infty$.
%244F

\spheader 244Xi Misalkan $\mu$ ukuran pencacahan pada $X=\{1,2,3\}$,
sehingga $\eusm L^0(\mu)=\BbbR^3$ dan $L^p(\mu)=L^0(\mu)$ dapat
diidentifikasi dengan $\BbbR^3$ untuk setiap $p\in[1,\infty]$. Uraikan
bola satuan $\{u:\|u\|_p\le 1\}$ dalam $\BbbR^3$ untuk $p=1$, $2$ dan
$\infty$.
%244F

\spheader 244Xj Pada titik mana sajakah argumen 244Hb gagal jika kita
mencoba menerapkannya pada $L^{\infty}$ dengan $\|\,\|_{\infty}$?
%244H

\spheader 244Xk Misalkan $p\in\coint{1,\infty}$. (i) Tunjukkan bahwa
$|a^p-b^p|\ge|a-b|^p$ untuk semua $a$, $b\ge 0$.
\Hint{untuk $a>b$, diferensiasikan kedua ruas terhadap $a$.}
(ii) Misalkan $(X,\Sigma,\mu)$ suatu ruang ukur dan $U$ suatu subruang
linear dari $L^0(\mu)$ sedemikian sehingga ($\alpha$) $|u|\in U$ untuk
setiap $u\in U$ ($\beta$) $u^{1/p}\in U$ untuk setiap $u\in U$
($\gamma$) $U\cap L^1$ padat dalam $L^1=L^1(\mu)$. Tunjukkan bahwa
$U\cap L^p$ padat dalam $L^p=L^p(\mu)$. \Hint{periksa terlebih dahulu
bahwa $\{u:u\in U\cap L^1$, $u\ge 0\}$ padat dalam
$\{u:u\in L^1$, $u\ge 0\}$.} (iii) Gunakan ini untuk membuktikan 244H
dari 242M dan 242O.
%244H

\spheader 244Xl Untuk sebarang ruang ukur $(X,\Sigma,\mu)$, tuliskan
$M^{1,\infty}=M^{1,\infty}(\mu)$ untuk
$\{v+w:v\in L^1(\mu),\,w\in L^{\infty}(\mu)\}\subseteq L^0(\mu)$.
Tunjukkan bahwa $M^{1,\infty}$ merupakan subruang linear dari $L^0$
yang memuat $L^p$ untuk setiap $p\in[1,\infty]$, dan bahwa jika
$u\in L^0$, $v\in M^{1,\infty}$ dan $|u|\le |v|$, maka
$u\in M^{1,\infty}$. \Hint{$u=v\times w$ dengan
$|w|\le\chi X^{\ssbullet}$.}
%244M

\spheader 244Xm Misalkan $(X,\Sigma,\mu)$ dan $(Y,\Tau,\nu)$ dua ruang
ukur, dan misalkan $\Cal T^+$ himpunan operator linear
$T:M^{1,\infty}(\mu)\to M^{1,\infty}(\nu)$ sedemikian sehingga
($\alpha$) $Tu\ge 0$ setiap kali $u\ge 0$ dalam $M^{1,\infty}(\mu)$
($\beta$) $Tu\in L^1(\nu)$ dan $\|Tu\|_1\le\|u\|_1$ setiap kali
$u\in L^1(\mu)$ ($\gamma$) $Tu\in L^{\infty}(\nu)$ dan
$\|Tu\|_{\infty}\le\|u\|_{\infty}$ setiap kali
$u\in L^{\infty}(\mu)$. (i) Tunjukkan bahwa jika
$\phi:\Bbb R\to\Bbb R$ suatu fungsi konveks sedemikian sehingga
$\phi(0)=0$, dan $u\in M^{1,\infty}(\mu)$ sedemikian sehingga
$\bar\phi(u)\in M^{1,\infty}(\mu)$ (dengan menafsirkan
$\bar\phi:L^0(\mu)\to L^0(\mu)$ seperti dalam 241I), maka
$\bar\phi(Tu)\in M^{1,\infty}(\nu)$ dan
$\bar\phi(Tu)\le T(\bar\phi(u))$ untuk setiap $T\in\Cal T^+$.
(ii) Oleh karena itu, tunjukkan bahwa jika $p\in[1,\infty]$ dan
$u\in L^p(\mu)$, maka $Tu\in L^p(\nu)$ dan
$\|Tu\|_p\le\|u\|_p$ untuk setiap $T\in\Cal T^+$.
%244M

\sqheader 244Xn Misalkan $X$ sebarang himpunan, dan misalkan $\mu$ ukuran
pencacahan pada $X$. Dalam kasus ini lazim dituliskan $\ell^p(X)$ untuk
${\eusm L}^p(\mu)$, dan ruang itu diidentifikasi dengan $L^p(\mu)$.
Secara khusus, $L^2(\mu)$ menjadi diidentifikasi dengan $\ell^2(X)$,
ruang fungsi yang dapat dijumlahkan-kuadrat pada $X$. Tuliskan pernyataan
dan bukti hasil-hasil bagian ini yang disesuaikan dengan kasus khusus ini.
%244N

\spheader 244Xo Misalkan $(X,\Sigma,\mu)$ dan $(Y,\Tau,\nu)$ ruang-ruang
ukur dan $\phi:X\to Y$ suatu fungsi \imp. Tunjukkan bahwa peta
$g\mapsto g\phi:\eusm L^0(\nu)\to\eusm L^0(\mu)$ (241Xg) menginduksi
suatu peta yang mempertahankan norma dari $L^p(\nu)$ ke $L^p(\mu)$ untuk
setiap $p\in[1,\infty]$, dan juga suatu peta dari $M^{1,\infty}(\nu)$ ke
$M^{1,\infty}(\mu)$ yang termasuk dalam kelas $\Cal T^+$ dari 244Xm.
%244N

\leader{244Y}{Latihan lanjutan (a)}
%\spheader 244Ya
Misalkan $(X,\Sigma,\mu)$ suatu ruang ukur, dan
$(X,\tilde\Sigma,\tilde\mu)$ versi c.l.d.-nya. Tunjukkan bahwa
$\eusm L^p(\mu)\subseteq\eusm L^p(\tilde\mu)$ dan bahwa pembenaman ini
menginduksi suatu isomorfisme kisi Banach antara $L^p(\mu)$ dan
$L^p(\tilde\mu)$, untuk setiap $p\in\coint{1,\infty}$.
%244A

\spheader 244Yb Misalkan $(X,\Sigma,\mu)$ sebarang ruang ukur, dan
$p\in\coint{1,\infty}$. Tunjukkan bahwa $L^p(\mu)$ mempunyai sifat
supremum terhitung dalam arti 241Ye. \Hint{242Yh.}
%244A

\spheader 244Yc\dvAnew{2009.} Andaikan $(X,\Sigma,\mu)$ suatu ruang ukur,
dan bahwa $p\in\ooint{0,1}$, $q<0$ sedemikian sehingga
$\bover1p+\bover1q=1$. (i) Tunjukkan bahwa
$ab\ge\bover1pa^p+\bover1qb^q$ untuk semua bilangan real $a\ge 0$,
$b>0$. \Hint{tetapkan $p'=\bover1p$, $q'=\bover{p'}{p'-1}$,
$c=(ab)^p$, $d=b^{-p}$ dan terapkan 244Ea.}
(ii) Tunjukkan bahwa jika $f$, $g\in\eusm L^0(\mu)$ nonnegatif dan
$E=\{x:x\in\dom g$, $g(x)>0\}$, maka

\Centerline{$(\int_Ef^p)^{1/p}(\int_Eg^q)^{1/q}
\le\int f\times g$.}

\noindent(iii) Tunjukkan bahwa jika $f$, $g\in\eusm L^0(\mu)$
nonnegatif, maka

\Centerline{$(\int f^p)^{1/p}+(\int g^p)^{1/p}
\le(\int(f+g)^p)^{1/p}$.}
%244E

\spheader 244Yd\dvAformerly{2{}44Yc.}
Misalkan $(X,\Sigma,\mu)$ suatu ruang ukur, dan $Y$ subhimpunan dari $X$;
tuliskan $\mu_Y$ untuk ukuran subruang pada $Y$. Tunjukkan bahwa peta
kanonik $T$ dari $L^0(\mu)$ pada $L^0(\mu_Y)$ (241Yg) memuat suatu
surjeksi dari $L^p(\mu)$ pada $L^p(\mu_Y)$ untuk setiap
$p\in[1,\infty]$, dan juga suatu peta dari $M^{1,\infty}(\mu)$ ke
$M^{1,\infty}(\mu_Y)$ yang termasuk dalam kelas $\Cal T^+$ dari 244Xm.
Tunjukkan bahwa pernyataan berikut ekuivalen: (i) terdapat suatu
$p\in\coint{1,\infty}$ sedemikian sehingga $T\restr L^p(\mu)$ injektif;
(ii) $T:L^p(\mu)\to L^p(\mu_Y)$ mempertahankan norma untuk setiap
$p\in\coint{1,\infty}$; (iii) $F\cap Y\ne\emptyset$ setiap kali
$F\in\Sigma$ dan $0<\mu F<\infty$.
%244F

\spheader 244Ye\dvAformerly{2{}44Yd.}
Misalkan $(X,\Sigma,\mu)$ sebarang ruang ukur, dan
$p\in\coint{1,\infty}$. Tunjukkan bahwa norma $\|\,\|_p$ pada
$L^p(\mu)$ kontinu-urutan dalam arti 242Yg.
%244F

\spheader 244Yf\dvAformerly{2{}44Ye.}
Misalkan $(X,\Sigma,\mu)$ sebarang ruang ukur, dan $p\in[1,\infty]$.
Tunjukkan bahwa jika $A\subseteq L^p(\mu)$ terarah ke atas dan terbatas
norma, maka himpunan itu terbatas di atas. \Hint{242Yf.}
%244F

\spheader 244Yg\dvAformerly{2{}44Yf.}
Misalkan $(X,\Sigma,\mu)$ sebarang ruang ukur, dan $p\in[1,\infty]$.
Tunjukkan bahwa jika suatu himpunan tak kosong $A\subseteq L^p(\mu)$
terarah ke atas dan mempunyai supremum dalam $L^p(\mu)$, maka
$\|\sup A\|_p\le\sup_{u\in A}\|u\|_p$. \Hint{tinjau terlebih dahulu
kasus $0\in A$.}
%244F

\spheader 244Yh\dvAformerly{2{}44Yg.}
Misalkan $(X,\Sigma,\mu)$ suatu ruang ukur dan $u\in L^0(\mu)$.
(i) Tunjukkan bahwa
$I=\{p:p\in\coint{1,\infty},\,u\in L^p(\mu)\}$ merupakan suatu
interval. Berikan contoh untuk menunjukkan bahwa interval itu dapat
terbuka, tertutup, atau setengah terbuka. (ii) Tunjukkan bahwa
$p\mapsto p\ln\|u\|_p:I\to\Bbb R$ konveks. \Hint{jika $p<q$ dan
$t\in\ooint{0,1}$, perhatikan bahwa
$\int|u|^{tp+(1-t)q}\le\||u|^{pt}\|_{1/t}\||u|^{q(1-t)}\|_{1/(1-t)}$.}
(iii) Tunjukkan bahwa jika $p\le q\le r$ dalam $I$, maka
$\|u\|_q\le\max(\|u\|_p,\|u\|_r)$.
%244F mt24bits

\spheader 244Yi\dvAformerly{2{}44Yh.}
Misalkan $[a,b]$ suatu interval tertutup tak trivial dalam $\Bbb R$ dan
$F:[a,b]\to\Bbb R$ suatu fungsi; ambil $p\in\ooint{1,\infty}$.
Tunjukkan bahwa pernyataan berikut ekuivalen: (i) $F$ kontinu mutlak dan
turunannya $F'$ termasuk dalam $\eusm L^p(\mu)$, dengan $\mu$ ukuran
Lebesgue pada $[a,b]$ (ii)

\Centerline{$c=\sup\{\sum_{i=1}^n\Bover{|F(a_i)-F(a_{i-1})|^p}
{(a_i-a_{i-1})^{p-1}}:a\le a_0<a_1<\ldots<a_n\le b\}$}

\noindent berhingga, dan bahwa dalam hal ini $c=\|F'\|_p$.
\Hint{(i) jika $F$ kontinu mutlak dan $F'\in\eusm L^p$, gunakan
pertidaksamaan H\"older untuk menunjukkan bahwa
$|F(b')-F(a')|^p\le(b'-a')^{p-1}\int_{a'}^{b'}|F'|^p$ setiap kali
$a\le a'\le b'\le b$. (ii) Jika $F$ memenuhi syarat tersebut, tunjukkan
bahwa
$(\sum_{i=0}^n|F(b_i)-F(a_i)|)^p\le c(\sum_{i=0}^n(b_i-a_i))^{p-1}$
setiap kali $a\le a_0\le b_0\le a_1\le\ldots\le b_n\le b$, sehingga
$F$ kontinu mutlak. Ambil suatu barisan $\sequencen{F_n}$ dari fungsi
poligonal yang menghampiri $F$; gunakan 223Xj untuk menunjukkan bahwa
$F'_n\to F'$ hampir di mana-mana, sehingga
$\int|F'|^p\le\sup_{n\in\Bbb N}\int|F'_n|^p\le c^p$.}
%244F

\spheader 244Yj\dvAformerly{2{}44Yi.}
Misalkan $G$ suatu himpunan terbuka dalam $\BbbR^r$ dan tuliskan $\mu$
untuk ukuran Lebesgue pada $G$. Misalkan $C_k(G)$ himpunan fungsi kontinu
$f:G\to\Bbb R$ sedemikian sehingga
$\inf\{\|x-y\|:x\in G,\,f(x)\ne 0,\,y\in\BbbR^r\setminus G\}>0$
(dengan menghitung $\inf\emptyset$ sebagai $\infty$).
Tunjukkan bahwa untuk sebarang $p\in\coint{1,\infty}$, himpunan
$\{f^{\ssbullet}:f\in C_k(G)\}$ merupakan subruang linear padat dari
$L^p(\mu)$.
%244H

\spheader 244Yk\dvAformerly{2{}44Yj.}
Misalkan $U$ sebarang ruang Hilbert. (i) Tunjukkan bahwa jika
$C\subseteq U$ konveks (yakni, $tu+(1-t)v\in C$ setiap kali
$u$, $v\in C$ dan $t\in[0,1]$; lihat 233Xd), tertutup dan tak kosong,
serta $u\in U$, maka terdapat tepat satu $v\in C$ sedemikian sehingga
$\|u-v\|=\inf_{w\in C}\|u-w\|$, dan $\innerprod{u-v}{v-w}\ge 0$ untuk
setiap $w\in C$. (ii) Tunjukkan bahwa jika $h\in U^*$, terdapat tepat
satu $v\in U$ sedemikian sehingga $h(w)=\innerprod{w}{v}$ untuk setiap
$w\in U$. \Hint{terapkan (i) dengan $C=\{w:h(w)=1\}$, $u=0$.}
(iii) Tunjukkan bahwa jika $V\subseteq U$ suatu subruang linear tertutup,
maka terdapat tepat satu proyeksi linear $P$ pada $U$ sedemikian sehingga
$P[U]=V$ dan $\innerprod{u-Pu}{v}=0$ untuk semua $u\in U$, $v\in V$
($P$ bersifat `ortogonal'). \Hint{ambil $Pu$ sebagai titik $V$ yang
terdekat dengan $u$.}
%244N

\spheader 244Yl\dvAformerly{2{}44Yk.}
Misalkan $(X,\Sigma,\mu)$ suatu ruang probabilitas, dan $\Tau$ suatu
subaljabar-$\sigma$ dari $\Sigma$. Gunakan bagian (iii) dari 244Yk untuk
menunjukkan bahwa terdapat suatu proyeksi ortogonal
$P:L^2(\mu)\to L^2(\mu\restrp\Tau)$ sedemikian sehingga
$\int_FPu=\int_Fu$ untuk setiap $u\in L^2(\mu)$ dan $F\in\Tau$.
Tunjukkan bahwa $Pu\ge 0$ setiap kali $u\ge 0$ dan bahwa
$\int Pu=\int u$ untuk setiap $u$, sehingga $P$ mempunyai perluasan unik
menjadi suatu operator kontinu dari $L^1(\mu)$ pada
$L^1(\mu\restrp\Tau)$. Gunakan ini untuk mengembangkan teori ekspektasi
bersyarat tanpa menggunakan teorema Radon-Nikod\'ym.
%244N

\spheader 244Ym ({\smc Roselli \& Willem 02})\dvAnew{2009}
(i) Tetapkan $C=\coint{0,\infty}^2\subseteq\BbbR^2$. Misalkan
$\phi:C\to\Bbb R$ suatu fungsi kontinu sedemikian sehingga
$\phi(tz)=t\phi(z)$ untuk semua $z\in C$.
Tunjukkan bahwa $\phi$ konveks (definisi: 233Xd) jika dan hanya jika
$t\mapsto\phi(1,t):\coint{0,\infty}\to\Bbb R$ konveks.
(ii) Tunjukkan bahwa jika $p\in\ooint{1,\infty}$ dan
$q=\bover{p}{p-1}$, maka $(s,t)\mapsto -s^{1/p}t^{1/q}$,
$(s,t)\mapsto -(s^{1/p}+t^{1/p})^p:C\to\Bbb R$ konveks.
(iii) Tunjukkan bahwa jika $p\in[1,2]$, maka
$(s,t)\mapsto|s^{1/p}+t^{1/p}|^p+|s^{1/p}-t^{1/p}|^p$ konveks.
(iv) Tunjukkan bahwa jika $p\in\coint{2,\infty}$, maka
$(s,t)\mapsto -|s^{1/p}+t^{1/p}|^p-|s^{1/p}-t^{1/p}|^p$ konveks.
(v) Gunakan (ii) dan 233Yj untuk membuktikan pertidaksamaan H\"older dan
Minkowski. (vi) Gunakan (iii) dan (iv) untuk membuktikan pertidaksamaan
Hanner. (vii) Sesuaikan metode tersebut untuk menjawab (ii) dan (iii)
dari 244Yc.
%244O

\spheader 244Yn\dvAnew{2009}(i) Tunjukkan bahwa setiap ruang hasil kali
dalam bersifat konveks seragam.
(ii) Misalkan $U$ suatu ruang Banach konveks seragam,
$C\subseteq U$ suatu himpunan konveks tertutup tak kosong, dan $u\in U$.
Tunjukkan bahwa terdapat tepat satu $v_0\in C$ sedemikian sehingga
$\|u-v_0\|=\inf_{v\in C}\|u-v\|$.
(iii) Misalkan $U$ suatu ruang Banach konveks seragam, dan
$A\subseteq U$ suatu himpunan terbatas tak kosong. Tetapkan
$\delta_0=\inf\{\delta:$ terdapat $u\in U$ sedemikian sehingga
$A\subseteq B(u,\delta)=\{v:\|v-u\|\le\delta\}\}$. Tunjukkan bahwa
terdapat tepat satu $u_0\in U$ sedemikian sehingga
$A\subseteq B(u_0,\delta_0)$.
%244O

\spheader 244Yo\dvAnew{2009} Misalkan $(X,\Sigma,\mu)$ suatu ruang ukur,
dan $u\in L^0(\mu)$. Andaikan $\sequencen{p_n}$ suatu barisan dalam
$[1,\infty]$ dengan limit $p\in[1,\infty]$. Tunjukkan bahwa jika
$\limsup_{n\to\infty}\|u\|_{p_n}$ berhingga, maka
$\lim_{n\to\infty}\|u\|_{p_n}$ terdefinisi dan sama dengan $\|u\|_p$.
%244O
}%end of exercises

\endnotes{
\Notesheader{244} Pada titik ini saya merasa kita harus meninggalkan
penyelidikan ruang fungsi lebih lanjut. Tahap berikutnya harus berupa
analisis abstrak yang sistematis atas kisi Banach umum. Ruang-ruang $L^p$
memberikan landasan kukuh bagi analisis semacam itu, karena ruang-ruang
tersebut memperkenalkan tema dasar kelengkapan norma, kelengkapan urutan,
dan identifikasi ruang dual. Dalam latihan-latihan saya telah mencoba
menunjukkan pentingnya lapisan konsep berikutnya: kekontinuan-urutan norma
dan hubungan antara keterbatasan norma dan keterbatasan urutan. Hal yang
belum sempat saya bahas adalah operator linear yang mempertahankan urutan
di antara ruang-ruang Riesz, yang menjadi kunci untuk memahami struktur
urutan ruang-ruang dual di sini. (Namun Anda dapat memulai dengan membaca
ulang teori operator ekspektasi bersyarat dalam 242J-242L,
%242J 242K 242L
243J dan 244M.) Semua topik ini dibahas dalam {\smc Fremlin 74} dan dalam
Bab 35 dan 36 pada jilid berikutnya.

Saya ingat salah seorang guru analisis saya dahulu mengatakan bahwa
ruang-ruang $L^p$ (untuk $p\ne 1$, $2$, $\infty$) entah bagaimana telah
masuk ke silabus dan tidak pernah dikeluarkan lagi. Saya sendiri akan
menyebutnya klasik, dalam arti bahwa ruang-ruang tersebut telah menjadi
bagian dari pengalaman bersama semua analis fungsional sejak analisis
fungsional dimulai; dan walaupun Anda bebas tidak menyukainya, Anda tidak
dapat mengabaikannya sebagaimana Anda tidak dapat mengabaikan Milton jika
mempelajari puisi Inggris. Pertidaksamaan H\"older, khususnya, mempunyai
banyak sekali penerapan; bukan hanya 244F dan 244K, melainkan juga,
misalnya, 244Xc-244Xd dan 244Yh-244Yi.

Ruang-ruang $L^p$, untuk $1\le p\le \infty$, membentuk semacam kontinum.
Ditinjau dari konsep-konsep yang dibahas di sini, tidak ada perbedaan yang
dapat ditarik di antara berbagai ruang $L^p$ untuk $1<p<\infty$, kecuali
pengamatan bahwa norma $L^2$ merupakan norma hasil kali dalam, yang
bersesuaian dengan geometri Euklides pada subruang-subruang berdimensi
hingganya. Untuk membedakan ruang-ruang $L^p$ yang lain, kita memerlukan
konsep-konsep yang jauh lebih halus dalam geometri ruang bernorma.

Ditinjau dari teorema-teorema yang diberikan di sini, $L^1$ tampaknya
lebih dekat dengan rentang tengah $L^p$ untuk $1<p<\infty$ daripada
$L^{\infty}$; jadi, untuk semua $1\le p<\infty$, $L^p$ lengkap Dedekind
(tidak bergantung pada ruang ukur yang terlibat), ruang $S$ dari kelas
ekuivalensi fungsi sederhana padat dalam $L^p$ (sekali lagi, untuk setiap
ruang ukur), dan dual $(L^p)^*$ (hampir) dapat diidentifikasi sebagai
ruang fungsi yang lain. Semua ini patut dipandang, dengan satu atau lain
cara, sebagai akibat kekontinuan-urutan norma $L^p$ untuk $p<\infty$.
Hambatan utama terhadap identifikasi universal $(L^1)^*$ dengan
$L^{\infty}$ adalah bahwa untuk ruang ukur yang bukan $\sigma$-hingga,
ruang $L^{\infty}$ dapat tidak memadai, dan bukan karena adanya patologi
dalam ruang $L^1$ itu sendiri. (Setidaknya pokok ini hendak saya bahas
kembali dalam Volume 3.) Ada pula pokok bahwa untuk ruang ukur yang tidak
semihingga, himpunan yang murni tak hingga dapat menyumbang pada
$L^{\infty}$ tanpa sumbangan yang bersesuaian pada $L^1$. Untuk
$1<p<\infty$, kedua masalah ini tidak dapat timbul. Setiap anggota dari
setiap $L^p$ demikian didukung sepenuhnya oleh bagian $\sigma$-hingga
dari ruang ukur, dan hal yang sama berlaku bagi dualnya -- lihat bagian
(c) dari bukti 244K.

Tentu saja $L^1$ memang mempunyai geometri yang nyata berbeda dari
ruang-ruang $L^p$ lainnya. Tanda pertama ialah bahwa ia tidak refleksif
sebagai ruang Banach (kecuali ketika berdimensi hingga), sedangkan untuk
$1<p<\infty$ identifikasi $(L^p)^*$ dengan $L^q$ dan $(L^q)^*$ dengan
$L^p$, dengan $q=p/(p-1)$, menunjukkan bahwa pembenaman kanonik $L^p$
ke dalam $(L^p)^{**}$ surjektif, yaitu bahwa $L^p$ refleksif. Namun bahkan
ketika $L^1$ berdimensi hingga, bola satuan $L^1$ dan $L^{\infty}$ jelas
berbeda jenis dari bola satuan $L^p$ untuk $1<p<\infty$; bola-bola itu
mempunyai sudut, alih-alih membulat dengan mulus (244Xh-244Xi).
Ungkapan langsung dari perbedaan tersebut terdapat dalam 244O. Seperti
biasa, kasus $p=2$ jauh lebih penting daripada kasus umum dan luar biasa
lebih mudah (244Yn(i)); dan perhatikan bagaimana pertidaksamaan Hanner
berbalik pada $p=2$. (Lihat 244Yc untuk pembalikan pertidaksamaan
H\"older dan Minkowski pada $p=1$.) Ada keadaan ketika berguna untuk
mengetahui bahwa $\|\,\|_1$ dan $\|\,\|_{\infty}$ dapat dihampiri,
dengan cara yang dapat diuraikan secara tepat, oleh norma konveks seragam
(244Yo). Saya telah menuliskan bukti 244O berdasarkan kecerdikan dan
kalkulus lanjut, seperti bukti 244E. Dengan sedikit lebih banyak tentang
himpunan dan fungsi konveks, yang digambarkan dalam 233Yf-233Yj,
%233Yf 233Yg 233Yh 233Yi 233Yj
terdapat bukti alternatif yang mencolok (244Ym). Tentu saja bukti 244Ea
juga menggunakan kekonveksan, dalam arah yang terbalik.

Bukti 244K, yang mengidentifikasi $(L^p)^*$, merupakan perjalanan yang
cukup panjang, dan wajar untuk bertanya apakah kita benar-benar harus
bekerja sekeras itu, khususnya karena dalam kasus $L^2$ kita mempunyai
argumen yang jauh lebih mudah (244Yk). Tentu saja kita dapat bergerak
lebih cepat jika mengetahui sedikit lebih banyak tentang kisi Banach
(\S369 dalam Volume 3 memuat fakta-fakta yang relevan), meskipun rute ini
menggunakan beberapa teorema yang sama sulitnya dengan 244K sebagaimana
diberikan. Terdapat rute alternatif yang menggunakan geometri ruang
$L^p$, mengikuti gagasan 244Yk, tetapi saya tidak menganggapnya lebih
mudah, dan argumen yang saya sajikan di sini setidaknya mempunyai
keutamaan menggunakan sebagian gagasan yang sama dengan identifikasi
$(L^1)^*$ dalam 243G. Perbedaannya ialah bahwa sementara dalam 243G kita
mungkin harus merangkai suatu keluarga besar fungsi $g_F$ (bagian (b-v)
dari bukti), dalam 244K hanya terdapat fungsi-fungsi $g_n$ yang
terhitung; akibatnya kita dapat membuat argumen ini berlaku untuk ruang
ukur sembarang, bukan hanya ruang yang dapat dilokalkan.

Geometri ruang Hilbert memberi kita suatu pendekatan terhadap ekspektasi
bersyarat yang tidak bergantung pada teorema Radon-Nikod\'ym (244Yl).
Namun, untuk mengubah gagasan-gagasan ini menjadi bukti teorema
Radon-Nikod\'ym itu sendiri diperlukan daya penentuan dan kecerdikan yang
dapat digunakan dengan lebih baik di tempat lain.

Argumen kekonveksan 233J/242K dapat digunakan pada banyak operator selain
ekspektasi bersyarat (lihat 244Xm). Kelas `$\Cal T^+$' yang diuraikan di
sana sebenarnya bukan kelas terbesar yang membuat argumen ini berlaku;
saya membawa gagasan-gagasan tersebut lebih jauh dalam Bab 37. Masih ada
jauh lebih banyak yang dapat dikatakan jika Anda menempatkan sebarang
pasangan ruang $L^p$ sebagai pengganti $L^1$ dan $L^{\infty}$ dalam
244Xl. 244Yh merupakan permulaan, tetapi untuk hasil sesungguhnya
(`teorema kekonveksan Riesz') saya merujuk Anda ke {\smc Zygmund 59},
XII.1.11 atau {\smc Dunford \& Schwartz 57}, VI.10.11.
}%end of notes


\discrpage
